\subsection{The Krein-Rutman theorem}

Lemma 3. --- Let \((a_n)_{n \geqslant 0}\) be a sequence of positive real numbers such that the power series

\[
f (z) = \sum_ {n \geqslant 0} a _ {n} z ^ {n}
\]

has a finite radius of convergence r \textgreater{} 0. Let D ⊂ C be the open disk with center 0 and radius r. There does not exist a holomorphic function \(\widetilde{f}\) defined on an open neighborhood U of r in C that coincides with f on U ∩ D.

Suppose that such a holomorphic function \(\tilde{f}\), defined on an open neighborhood U of \(r\), exists. There exist real numbers \(s < r\) and \(\delta > 0\) such that \(s + \delta > r\) and such that the open disk D' with center \(s\) and radius \(\delta\) is contained in U. The power series expansion of \(\tilde{f}\) at the point \(s\) (VAR, R1, p. 26, 3.2.1) then converges for every \(z\) in the disk with center 0 and radius \(\delta\) (VAR, R1, p. 29, 3.3.4). This series is

\[
\widetilde {f} _ {s} (z) = \sum_ {n = 0} ^ {+ \infty} \frac {\widetilde {f} ^ {(n)} (s)}{n !} z ^ {n} = \sum_ {n = 0} ^ {+ \infty} \frac {f ^ {(n)} (s)}{n !} z ^ {n}
\]

(VAR, R1, p. 27, 3.2.4). Since \(s \in D\) , we have

\[
f ^ {(n)} (s) = \sum_ {k = n} ^ {+ \infty} k (k - 1) \dots (k - n + 1) a _ {k} s ^ {k - n}
\]

(VAR, R1, p. 28, 3.2.11). Take \(z\) such that \(0 < z < \delta\) . Since \(a_{k} \geqslant 0\) , we have

\[
\begin{array}{r l} & {\widetilde {f} _ {s} (z) = \sum_ {n = 0} ^ {+ \infty} \biggl (\sum_ {k = n} ^ {+ \infty} \frac {k (k - 1) \cdots (k - n + 1)}{n !} a _ {k} s ^ {k - n} \biggr) z ^ {n}} \\ & {\qquad = \sum_ {k = 0} ^ {+ \infty} \biggl (a _ {k} \sum_ {n = 0} ^ {k} \frac {k !}{n ! (k - n) !} s ^ {k - n} z ^ {n} \biggr) = \sum_ {k = 0} ^ {+ \infty} a _ {k} (s + z) ^ {k}} \end{array}
\]

(TG, III, p. 40). When \(s + z > r\) , the convergence of this expression contradicts the hypothesis that the radius of convergence of the power series \(f\) is equal to \(r\) .

Let E be a real Banach space. Let C be a pointed convex cone in E. The vector space generated by C is equal to C -- C (EVT, II, p. 12, cor. 1). In particular, the cone C is total in E if and only if C -- C is dense in E. Recall that C is said to be salient if C ∩ (−C) is reduced to 0 (EVT, II, p. 11).

The polar C° of C is the set of continuous linear forms \(\ell \in E'\) such that \(\ell(x) \geqslant 0\) for all \(x \in C\) (EVT, II, p. 47, prop. 4). By the bipolar theorem (EVT, II, p. 48, th. 1), if C is a closed pointed convex cone, then C°° = C.

Lemma 4. --- Let E be a real Banach space and \(u \in \mathcal{L}(\mathrm{E}_{(\mathbf{C})})\) a nonzero endomorphism of \(\mathrm{E}_{(\mathbf{C})}\) . Let C be a total convex cone in E. There exists \(x \in \mathrm{C}\) such that \(u(x) \neq 0\) .

Indeed, if \(u\) vanishes on C, then \(u\) vanishes on C - C, hence on E, and therefore on \(\mathrm{E}_{(\mathbf{C})}\) .

THEOREM 4 (Krein--Rutman). --- Let E be a complete normable space over R. Let C be a closed total salient convex cone in E and let \(u \in \mathcal{L}(\mathrm{E})\) be a compact linear map such that \(u(\mathrm{C}) \subset \mathrm{C}\) . If the spectral radius \(\varrho(u)\) is \textgreater0, then \(\varrho(u)\) is an isolated point of the spectrum of u, and there exists a nonzero eigenvector \(x \in C\) of u for the eigenvalue \(\varrho(u)\) .

Let \(E_{(\mathbf{C})}\) be the complexification of E and \(u_{(\mathbf{C})}\) the endomorphism of \(E_{(\mathbf{C})}\) obtained by extension of scalars from u. The spectral radius of \(u_{(\mathbf{C})}\) is equal to \(\varrho(u)\) (I, p. 86); it is simply denoted by \(\varrho\) . Since \(\varrho > 0\) , the complex spectrum of \(u\) is not reduced to 0. There therefore exists \(\lambda_0 \in \mathrm{Sp}_\mathrm{s}(u_{(\mathbf{C})})\) such that \(|\lambda_0| = \varrho\) (prop. 5 of III, p. 90). Let \(e_0\) be the spectral projector of \(u_{(\mathbf{C})}\) associated with \(\lambda_0\) .

The resolvent \(\lambda \mapsto \mathrm{R}(u_{(\mathbf{C})},\lambda) = (\lambda -u_{(\mathbf{C})})^{-1}\) is holomorphic on the complement of the spectrum of \(u_{(\mathbf{C})}\) (th. 1 of I, p. 24). The complex number \(\lambda_0\) is a pole of the resolvent and its residue is the spectral projector \(e_0\) (I, p. 131).

Let \(y\) and \(y'\) be elements of E such that \(y + iy' \in \mathrm{E}_{(\mathbf{C})}\) is an eigenvector of \(u_{(\mathbf{C})}\) for the eigenvalue \(\lambda_0\) . Since \(e_0(y + iy') = y + iy' \neq 0\) , there exists an element \(x_0 \in \mathrm{C}\) such that \(e_0(x_0) \neq 0\) (lemma 4). Since C is closed and pointed, its polar \(\mathrm{C}^\circ\) is total (EVT, II, p. 48, cor. 1), and hence there exists a linear form \(\ell_0 \in \mathrm{C}^\circ\) such that \(\langle e_0(x_0), \ell_0 \rangle \neq 0\) .

Consider the function \(f\) defined by \(f(0)=0\) and by

\[
f (z) = \langle \mathrm{R} (u _ {(\mathbf {C})}, z ^ {- 1}) x _ {0}, \ell_ {0} \rangle
\]

for \(z \in \mathbf{C}\) such that \(z^{-1} \notin \mathrm{Sp}(u_{(\mathbf{C})})\) . This function satisfies

\[
f (z) = \ell_ {0} \bigg (\Bigl (\sum_ {n = 0} ^ {+ \infty} z ^ {n + 1} u _ {(\mathbf {C})} ^ {n} \Bigr) x _ {0} \bigg) = \sum_ {n = 0} ^ {\infty} \langle u ^ {n} (x _ {0}), \ell_ {0} \rangle z ^ {n + 1}\tag{1}
\]

for \(|z| < 1 / \varrho\) (Theorem 1 of I, p. 24, d)) and is therefore holomorphic in the disk centered at 0 with radius \(1 / \varrho\) .

There exists a holomorphic function \(\widetilde{\mathrm{R}}\) defined on an open neighborhood U of \(\lambda_0\) and with values in \(\mathcal{L}(\mathrm{E})\) such that for \(z\) belonging to \(\mathrm{U} - \{\lambda_0\}\) , we have

\[
\mathrm{R} (u _ {(\mathbf {C})}, z) = \widetilde {\mathrm{R}} (z) + \sum_ {n = 0} ^ {+ \infty} (z - \lambda_ {0}) ^ {- n - 1} (u _ {(\mathbf {C})} - \lambda_ {0}) ^ {n} e _ {0}
\]

(prop. 17 of I, p. 83). For \(z\) such that \(z^{-1} \in \mathrm{U}\) and \(z \neq 1 / \lambda_0\) , we therefore have

\[
f (z) = \langle \widetilde {\mathrm{R}} (z ^ {- 1}) x _ {0}, \ell_ {0} \rangle + \sum_ {n = 0} ^ {+ \infty} (z ^ {- 1} - \lambda_ {0}) ^ {- n - 1} \langle (u _ {(\mathbf {C})} - \lambda_ {0}) ^ {n} e _ {0} (x _ {0}), \ell_ {0} \rangle .
\]

The term of the series corresponding to \(n=0\) is \((z^{-1}-\lambda_{0})^{-1}\langle e_{0}(x_{0}),\ell_{0}\rangle\) . Since \(\langle e_{0}(x_{0}),\ell_{0}\rangle\neq0\) , the uniqueness of the Laurent expansion (VAR, R1, p. 30, 3.3.9) implies that \(f\) does not extend to a holomorphic function in a neighborhood of \(1/\lambda_{0}\) . In particular, the radius of convergence of the power series expansion (1) of the function \(f\) at the point 0 is equal to \(1/\varrho\) .

The coefficients of the power series (1) are \(\langle u^n(x_0),\ell_0\rangle\geqslant 0\) since \(u(\mathrm{C})\subset\mathrm{C}\) and \(\ell_{0}\in\mathrm{C}^{\circ}\) . By Lemma 3, the function \(f\) does not extend to a holomorphic function in a neighborhood of \(1/\varrho\) . The resolvent of \(u_{(\mathbf{C})}\) therefore cannot be extended to a holomorphic function in a neighborhood of \(\varrho\) , that is, \(\varrho\in\mathrm{Sp}(u)\) . This implies that \(\varrho\) is an eigenvalue of \(u_{(\mathbf{C})}\) (prop. 5 of III, p. 90). Since \(\varrho\) is real, it is also an eigenvalue of \(u\) .

Let \(d \geqslant 1\) be the order of the pole of the resolvent of \(u_{(\mathbf{C})}\) at \(\varrho\) . Let \(e\) be the endomorphism

\[
e = \lim _ {z \to \varrho} (z - \varrho) ^ {d} \mathrm{R} (u _ {(\mathbf {C})}, z).
\]

It is nonzero and commutes with \(u_{(\mathbf{C})}\) , and its image is contained in the eigenspace of \(u_{(\mathbf{C})}\) relative to \(\varrho\) (cf. no. 3 of I, p. 131). Now let \(\ell\) be an element of C° and \(x\) an element of C. By Theorem 1 of I, p. 24, d), we have

\[
\langle e(x),\ell \rangle = \lim_{\substack{z\to \varrho \\ z > \varrho}}(z - \varrho)^{d}\sum_{n\geqslant 0}\langle u^{n}(x),\ell \rangle z^{-n - 1}\geqslant 0.
\]

The bipolar theorem (EVT, II, p. 48, th. 1) implies that \(e(x) \in \mathrm{C}\) for all \(x \in \mathrm{C}\) . Since \(\mathrm{C}\) is total, there exists \(x \in \mathrm{C}\) such that \(e(x) \neq 0\) (Lemma 4), and then \(e(x)\) belongs to \(\mathrm{C}\) and is an eigenvector of \(u\) for the eigenvalue \(\varrho\) , as desired.

COROLLARY. --- Let E be a complete normable space over R. Let C be a closed, total, salient convex cone in E, and \(u \in \mathcal{L}(E)\) be a compact linear map such that \(u(\mathrm{C}) \subset \mathrm{C}\) . If the spectral radius \(\varrho(u)\) of u is \textgreater{} 0, then \(\varrho(^{t}u) = \varrho(u)\) is an eigenvalue of \(^{t}u\) , and \(^{t}u\) admits an eigenvector relative to \(\varrho(u)\) in \(C^{\circ}\) .

One has \(\varrho(^{t}u) = \varrho(u)\) (prop. 3 of I, p. 131). By corollary 1 of III, p. 9, the endomorphism \(^{t}u\) of E' is compact. Moreover, we have \(^{t}u(C^{\circ}) \subset C^{\circ}\) ; the assertion then follows from the Krein-Rutman theorem applied to \(^{t}u\) , and to the closed convex cone \(C^{\circ}\) since the latter is salient (because C is total) and total (because C is salient).

Remark. --- The hypothesis \(\varrho(u) > 0\) cannot be omitted in general in the Krein–Rutman theorem. For example, let V be the endomorphism of the Banach space \(\mathcal{C}([0,1])\) defined by

\[
\mathrm{V} (f) (x) = \int_ {0} ^ {x} f (y) d y
\]

for \(f \in \mathcal{C}([0,1])\) . The map V is compact and its spectral radius is zero (exercise 1 of I, p. 187). It preserves the closed, total, salient convex cone in \(\mathcal{C}([0,1])\) consisting of positive functions, and has no eigenvalue (loc. cit.).

The following proposition describes a sufficient condition for an endomorphism preserving a cone to have a strictly positive spectral radius, and it then refines the Krein-Rutman theorem.

PROPOSITION 8. --- Let E be a nonzero complete normable space over R. Let C be a closed salient convex cone with nonempty interior \(\mathring{\mathrm{C}}\) in E, and let \(u \in \mathcal{L}(\mathrm{E})\) be a compact linear map such that \(u(\mathrm{C} - \{0\}) \subset \mathring{\mathrm{C}}\) .

\begin{enumerate}
\def\labelenumi{\alph{enumi})}
\item
  One has \(\varrho(u) > 0\) and there exists an eigenvector \(x_0\) of \(u\) in \(\mathring{C}\) for the eigenvalue \(\varrho(u)\) ;
\item
  The spectral projector of u corresponding to \(\varrho(u)\) has rank 1;
\item
  For every eigenvalue \(\lambda \neq \varrho(u)\) of \(u_{(\mathbf{C})}\) , one has \(|\lambda| < \varrho(u)\) ;
\item
  Let F be a subspace of E stable under u such that \((\mathrm{C}-\{0\})\cap\mathrm{F}\) is nonempty. Then \(x_{0}\in F\) . In particular, the only eigenvectors of u in C are the multiples of \(x_{0}\) .
\end{enumerate}

Let us prove two preliminary lemmas.

Lemma 5. --- Let E be a real Banach space, C a convex cone in E, and \(u \in \mathcal{L}(\mathrm{E})\) such that \(u(\mathrm{C} - \{0\}) \subset \mathring{\mathrm{C}}\) . Let \(\ell \in \mathrm{C}^{\circ}\) be an eigenvector of \(^{t}u\) . Then the kernel of \(\ell\) is stable under \(u\) and does not meet \(\mathrm{C} - \{0\}\) .

Let \(\lambda \in \mathbf{R}\) such that \(^t u(\ell) = \lambda \ell\) . For every \(x \in \mathrm{Ker}(\ell)\) , we have

\[
\langle u (x), \ell \rangle = \langle x, ^ {t} u (\ell) \rangle = \lambda \langle x, \ell \rangle = 0,
\]

hence \(\operatorname{Ker}(\ell)\) is stable under \(u\) .

Let \(x\) be a nonzero element of \(\operatorname{Ker}(\ell)\) . For every element \(y\) of E such that \(\langle y, \ell \rangle < 0\) , we have \(\langle u(x) + y, \ell \rangle < 0\) , whence it follows that \(u(x) + y \notin \mathrm{C}\) since \(\ell \in \mathrm{C}^\circ\) . We deduce that \(u(x)\) does not belong to \(\mathring{\mathrm{C}}\) . Since \(u(\mathrm{C} - \{0\}) \subset \mathring{\mathrm{C}}\) , this implies that \(x \notin \mathrm{C}\) .

Lemma 6. --- Let E be a finite-dimensional real vector space and let B be a compact convex neighborhood of 0 in E. Let u be an endomorphism of E such that \(u(B) \subset \mathring{B}\) . The complex spectrum of u is then contained in the unit disk of C and, in particular, it does not meet the circle with center 0 and radius 1 in C.

By replacing B with \(B \cap (-B)\) , we reduce to the case where B is balanced. The gauge p of B (EVT, II, p. 28, ex. 3) is then a norm on E, which defines its topology (EVT, I, p. 14, th. 2). The hypothesis implies that \(p(u(x)) < p(x)\) for every element x of E - \{0\}, hence that the spectral radius of u is \textless{} 1; since this is also the spectral radius of \(u_{(\mathbf{C})}\) (cf. I, p. 86), the conclusion follows.

Let us now prove the proposition.

Denote by \(x \preccurlyeq y\) the order relation on E associated with the convex cone C (EVT, II, p. 13, n° 5), that is, \(x \preccurlyeq y\) if and only if \(y - x \in \mathrm{C}\) . We have \(u(x) \preccurlyeq u(y)\) if \(x \preccurlyeq y\) . Since \(\mathring{\mathrm{C}}\) is nonempty, the vector space C -- C generated by C (EVT, II, p. 12, cor. 1) contains a neighborhood of 0, hence the cone C is total.

Let us prove assertion \(a\) ). Let \(\varrho = \varrho(u)\) . Let \(y_0\) be an element of \(\mathring{\mathrm{C}}\) . We have \(y_0 \neq 0\) . Let \(r > 0\) be such that the closed ball with center \(y_0\) and radius \(r\) is contained in C. For every \(y \in \mathrm{E} - \{0\}\) , we have \(y_0 - r \| y\|^{-1}y \in \mathrm{C}\) , hence \(y \preccurlyeq r^{-1}\| y\| y_0\) for all \(y \in \mathrm{E}\) .

Since \(y_{0} \neq 0\) , the hypotheses imply that \(u(y_{0}) \in \mathring{\mathrm{C}}\) . There therefore exists t \textgreater{} 0 such that \(tu(y_{0}) - y_{0} \in \mathrm{C}\) . Set v = tu. This is a compact endomorphism of E such that \(v(\mathrm{C}) \subset \mathrm{C}\) and \(v(y_{0}) \succcurlyeq y_{0}\) . For every integer \(n \geqslant 1\) , we have

\[
y _ {0} \preccurlyeq v ^ {n} (y _ {0}) \preccurlyeq r ^ {- 1} \| v ^ {n} (y _ {0}) \| y _ {0} \preccurlyeq r ^ {- 1} \| v ^ {n} \| \| y _ {0} \| y _ {0},
\]

hence \((r^{-1}\|v^{n}\|\|y_{0}\|-1)y_{0}\in\mathrm{C}\) . Since C is salient, this implies that \(t^{n}\|u^{n}\|=\|v^{n}\|\geqslant r/\|y_{0}\|\) , whence \(\varrho\geqslant t^{-1}>0\) (I, p. 20, prop. 1).

By the Krein--Rutman theorem (th. 4), the real number \(\varrho\) is an eigenvalue of \(u\) , and there exists an eigenvector \(x_0\) of \(u\) in C for the eigenvalue \(\varrho\) . We have \(\varrho x_0 = u(x_0) \in \mathring{\mathrm{C}}\) by hypothesis, hence \(x_0 \in \mathring{\mathrm{C}}\) . This establishes assertion \(a\) ).

Let K be the intersection of the spectrum of \(u_{(\mathbf{C})}\) and the circle with center 0 and radius \(\varrho\) in C. Since u is compact and \(\varrho > 0\) , the set K is finite and the image of the spectral projector \(e_{K}\) is a finite-dimensional subspace of \(\mathrm{E}_{(\mathbf{C})}\) (prop. 2 of III, p. 83 and prop. 5 of III, p. 90). Since K is stable under complex conjugation, the image of \(e_{K}\) is a finite-dimensional subspace of \(\mathrm{E}_{(\mathbf{C})}\) rational over R (A, V, p.60, prop. 6); denote by F the subspace of E such that \(\mathrm{F}_{(\mathbf{C})}\) is equal to the image of \(e_{K}\) . The space F is stable under u and contains the eigenspace of u relative to \(\varrho\) , hence F is nonzero.

To prove the assertions \(b)\) and \(c),\) it suffices to prove that \(F\) is of dimension 1.

Denote by \(v\) the endomorphism of F deduced from \(u\) by passing to subspaces. Let \(C_F = C \cap F\) . This is a closed salient convex cone with nonempty interior in F (since \(x_0 \in \mathring{C} \cap F\) ), hence total; it satisfies \(v(C_F - \{0\}) \subset \mathring{C}_F\) , and in particular is stable under \(v\) . Since \(\varrho(v) \leqslant \varrho\) and \(x_0 \in F\) , we have \(\varrho(v) = \varrho\) . By the corollary of Theorem 4, applied to \(C_F\) and \(v\) there exists an eigenvector \(\ell\) of \(^t v\) in \(C_F^\circ\) for the eigenvalue \(\varrho(v) = \varrho > 0\) .

The subspace \(\mathrm{Ker}(\ell)\) of F is stable under v. Let w be the endomorphism of \(\mathrm{Ker}(\ell)\) deduced from \(\varrho^{-1}v\) by passing to the quotient subspaces; its spectrum is contained in the unit circle. The set B of \(y \in \mathrm{Ker}(\ell)\) such that \(x_{0} + y \in C_{F}\) is a closed convex neighborhood of 0 in \(\mathrm{Ker}(\ell)\) . For every \(y \in B\) and every \(z \in \mathrm{Ker}(\ell)\) , we have

\[
x _ {0} + (w (y) + z) = \varrho^ {- 1} (v (x _ {0} + y) + \varrho z)
\]

which belongs to \(C_{F}\) if the norm of z is sufficiently small, hence \(w(y) \in \mathring{\mathrm{B}}\) . The set B is bounded: indeed, if there existed a sequence \((y_{n})_{n \in \mathbf{N}}\) in \(\operatorname{Ker}(\ell)\) such that \(\|y_{n}\| \to +\infty\) and \(x_{0} + y_{n} \in C_{F}\) then, after passing to a subsequence, we would have \(y_{n}/\|y_{n}\| \to y\) where \(y \in \operatorname{Ker}(\ell)\) is nonzero, and \(y = \lim \|y_{n}\|^{-1}(x_{0} + y_{n})\) would belong to \(C_{F}\) , contradicting Lemma 5. The set B is therefore compact.

One then deduces from Lemma 6, applied to \(\mathrm{Ker}(\ell)\) , to B, and to \(w\) , that the complex spectrum of \(w\) does not meet the circle with center 0 and radius 1; this implies that the spectrum of \(w\) is empty, which means that \(\mathrm{Ker}(\ell)\) is reduced to \(\{0\}\) , hence that F has dimension 1. The assertions \(b\) and \(c\) are therefore established.

The linear map \({}^t u\) is compact (Corollary 1 of III, p. 9) and \({}^t u(\mathrm{C}^\circ) \subset \mathrm{C}^\circ\) ; moreover \(\varrho(^t u) = \varrho > 0\) (Prop. 3 of I, p. 131). By the corollary of Theorem 4, there exists in \(\mathrm{C}^\circ\) an eigenvector \(\ell \neq 0\) of \({}^t u\) for the eigenvalue \(\varrho\) . The kernel of \(\ell\) is stable under \(u\) and does not meet \(\mathrm{C} - \{0\}\) (Lemma 5). In particular, one has \(x_0 \notin \operatorname{Ker}(\ell)\) .

Let F be a subspace of E stable under u. One has the decomposition \(\mathrm{F} = (\mathrm{F} \cap \mathbf{R}x_{0}) \oplus (\mathrm{F} \cap \operatorname{Ker}(\ell))\) . If F does not contain \(x_{0}\) , we therefore have \(\mathrm{F} \subset \operatorname{Ker}(\ell)\) , then F does not meet \(C - \{0\}\) . This establishes d) and concludes the proof of the proposition.

COROLLARY. --- Let E be a nonzero complete normable space over R. Let C be a closed convex salient cone with nonempty interior \(\dot{C}\) in E, and let u be a compact endomorphism of E such that \(u(\mathrm{C}) \subset \mathrm{C}\) . Suppose that there exists an integer \(k \geqslant 1\) such that \(u^{k}(\mathrm{C} - \{0\}) \subset \mathring{\mathrm{C}}\) .

\begin{enumerate}
\def\labelenumi{\alph{enumi})}
\item
  One has \(\varrho(u) > 0\) and there exists an eigenvector \(x_0\) of \(u\) in \(\dot{C}\) for the eigenvalue \(\varrho(u)\) ;
\item
  The spectral projector of u corresponding to \(\varrho(u)\) has rank 1;
\item
  For every eigenvalue \(\lambda \neq \varrho(u)\) of \(u_{(\mathbf{C})}\) , one has \(|\lambda| < \varrho(u)\) ;
\item
  The only eigenvectors of u in C are the multiples of \(x_0\) .
\end{enumerate}

Set \(\varrho = \varrho(u)\) . One can apply Th. 4 to \(u\) and Prop. 8 to \(u^k\) . There therefore exists an eigenvector \(x_0\) of \(u\) in C for the eigenvalue \(\varrho\) . Since \(0 < \varrho(u^k) = \varrho^k\) (formula (4) of I, p. 21), one has in particular \(\varrho > 0\) , and since \(x_0\) is an eigenvector of \(u^k\) for the eigenvalue \(\varrho(u^k)\) , we have \(x_0 \in \mathring{\mathrm{C}}\) .

One has \(\mathrm{Sp}(u_{(\mathbf{C})}^{k})=\mathrm{Sp}(u_{(\mathbf{C})})^{k}\) (Remark 4 of I, p. 2), hence every eigenvalue \(\lambda\in C\) of \(u_{(\mathbf{C})}\) such that \(|\lambda|=\varrho\) satisfies \(\lambda^{k}=\varrho^{k}\) , and since every eigenvector corresponding to \(\lambda\) is an eigenvector of \(u^{k}\) for \(\varrho^{k}\) , hence proportional to \(x_{0}\) , we have \(\lambda=\varrho\) .

If the spectral projector of u for the eigenvalue \(\varrho\) had rank at least 2, then so would that of \(u^{k}\) for the eigenvalue \(\varrho^{k}\) (cf. no. 2 of I, p. 129).

Finally, if \(x \in C\) is an eigenvector of u, it is also one for \(u^{k}\) , hence is proportional to \(x_{0}\) .

Let n be an integer \(\geqslant 1\) . The set \(R_{+}^{n}\) is a closed convex pointed salient cone in \(R^{n}\) , with nonempty interior equal to \((\mathbf{R}_{+}^{*})^{n}\) .

Let \(A = (a_{i,j})\) be a real matrix of type \((n,n)\) , and \(u_A\) the endomorphism \(x \mapsto Ax\) of \(\mathbf{R}^n\) . Let \(\varrho = \varrho(u_A)\) be its spectral radius.

Lemma 7. --- One has \(u_A(\mathbf{R}_+^n) \subset \mathbf{R}_+^n\) if and only if \(a_{i,j} \geqslant 0\) for all i and j, and \(u_A(\mathbf{R}_+^n - \{0\}) \subset (\mathbf{R}_+^*)^n\) if and only if \(a_{i,j} > 0\) for all i and j.

Let \((e_{1},\ldots,e_{n})\) be the canonical basis of \(R^{n}\) . The vectors \(e_{i}\) belong to \(R_{+}^{n}\) and \(u_{A}(e_{i})\in\mathbf{R}_{+}^{n}\) for all i (resp. \(u_{A}(e_{i})\in(\mathbf{R}_{+}^{*})^{n}\) for all i) if and only if \(a_{i,j}\geqslant0\) for all i and j (resp. \(a_{i,j}>0\) for all i and j). The result follows by linearity.

THEOREM 5 (Perron--Frobenius). --- a) If \(a_{i,j} \geqslant 0\) for all \(i\) and \(j\) in \(\{1,\ldots,n\}\) , then the real number \(\varrho\) is an eigenvalue of \(A\) ; b) If \(a_{i,j} > 0\) for all \(i\) and \(j\) in \(\{1,\ldots,n\}\) , then \(\varrho > 0\) and the primary space of \(u_A\) relative to \(\varrho\) , that is, the nullspace of \(A - \varrho 1_{\mathbf{R}^n}\) , has dimension 1 and is generated by a vector \(x_0 \in (\mathbf{R}_+^*)^n\) . There is no other eigenvalue of \(A\) admitting an eigenvector in \(\mathbf{R}_+^n\) , and all complex eigenvalues \(\lambda \neq \varrho\) of \(A\) satisfy \(|\lambda| < \varrho\) .

If \(\varrho = 0\) , assertion \(a\) ) is elementary because 0 is then an eigenvalue of \(A\) . If \(\varrho > 0\) , assertion \(a\) ) follows from Theorem 4 in view of Lemma 7. Assertion \(b\) ) is, for the same reason, a consequence of Proposition 8.

Remark. --- Suppose there exists an integer \(k \geqslant 1\) such that all coefficients of \(A^{k}\) are \textgreater0 (such a matrix is sometimes called primitive). Then, by the corollary of Prop. 8, the spectral radius \(\varrho\) is an eigenvalue of A, the primary space of \(u_{A}\) relative to \(\varrho\) has dimension 1, and is generated by a vector \(x_{0} \in (\mathbf{R}_{+}^{*})^{n}\) . Moreover, there is no other complex eigenvalue of A admitting an eigenvector in \(R_{+}^{n}\) and all complex eigenvalues \(\lambda \neq \varrho\) of A satisfy \(|\lambda| < \varrho\) .

\specialchapter{Exercises}

\exercisesection{}

In Exercises 1 to 5 exclusively, K denotes a complete non-Archimedean non-discrete valued field, whose valuation is denoted by \(x \mapsto |x|\) . We also denote by \(\mathrm{G} \subset \mathbf{R}_{+}^{*}\) the image of \(\mathrm{K}^{*}\) under the map \(x \mapsto |x|\) . The Banach spaces considered are Banach spaces over \(\mathrm{K}\) .

\begin{enumerate}
\def\labelenumi{\arabic{enumi})}
\tightlist
\item
  \begin{enumerate}
  \def\labelenumii{\alph{enumii})}
  \tightlist
  \item
    Let E be a Banach space over K. There exists a norm \(x \mapsto \|x\|\) on E, defining the topology of E, such that \(\|x\| \in \overline{G}\) for all \(x \in E\) . Such a norm will be called round. If the valuation of K is not discrete, then every norm defining the topology of E is round.
  \end{enumerate}
\end{enumerate}

\begin{enumerate}
\def\labelenumi{\alph{enumi})}
\setcounter{enumi}{1}
\item
  Let E and F be Banach spaces over K. We equip the space \(\mathcal{L}(\mathrm{E};\mathrm{F})\) with the norm \(\| u\| = \sup_{x\in \mathrm{E}- \{0\}}\| u(x)\| /\| x\|\) . The space \(\mathcal{L}(\mathrm{E};\mathrm{F})\) is a Banach space over K. If the norm on E is round, then \(\| u\| = \sup_{\| x\| \leqslant 1}\| u(x)\|\) . If the norm on F is also round, then the norm on \(\mathcal{L}(\mathrm{E};\mathrm{F})\) is round. We denote \(\mathrm{E}' = \mathcal{L}(\mathrm{E};\mathrm{K})\) and we say that \(\mathrm{E}'\) is the dual space of E.
\item
  Let I be a set, F a Banach space over K, and \(\mathcal{C}_{0}(\mathrm{I},\mathrm{F})\) the K-vector space of functions \(f\colon I\to F\) such that \(\|f(i)\|\) tends to 0 along the filter of complements of finite subsets of I. Equipped with the norm \(\|f\|=\sup_{i\in I}\|f(i)\|\) , the space \(\mathcal{C}_{0}(\mathrm{I},\mathrm{F})\) is a Banach space over K, whose norm is round if that of F is. We denote \(\mathcal{C}_{0}(\mathrm{I})=\mathcal{C}_{0}(\mathrm{I},\mathrm{K})\) .
\end{enumerate}

\(d\) ) Let I be a set. For every \(i\in \mathrm{I}\) , we denote by \(\varphi_{i}\in \mathcal{C}_{0}(\mathrm{I})\) the characteristic function of \(\{i\}\) . Let F be a Banach space over K. The map \(u \mapsto (u(\varphi_i))_{i \in \mathrm{I}}\) is an isomorphism from the Banach space \(\mathcal{L}(\mathcal{C}_0(\mathrm{I});\mathrm{F})\) to the Banach space \(\mathcal{B}(\mathrm{I};\mathrm{K})\) (EVT, I, p. 4, example; cf. EVT, I, p. 23, exercise 7).

\begin{enumerate}
\def\labelenumi{\alph{enumi})}
\setcounter{enumi}{4}
\item
  Assume that the valuation of K is discrete. Let E be a Banach space over K whose norm is round. There exists a set I such that E is isometrically isomorphic to the Banach space \(\mathcal{C}_{0}(\mathrm{I})\) (cf. EVT, I, p. 23, exercise 7).
\item
  Assume that the valuation of K is discrete. Let E be a Banach space over K equipped with a round norm and F a closed vector subspace of E. There exists a continuous projector of E whose image is F and which has norm \(\leqslant 1\) .
\end{enumerate}

\begin{enumerate}
\def\labelenumi{\arabic{enumi})}
\setcounter{enumi}{1}
\tightlist
\item
  Let E and F be Banach spaces over K. We say that \(u \in \mathcal{L}(\mathrm{E};\mathrm{F})\) is completely continuous if it belongs to the closure in \(\mathcal{L}(\mathrm{E};\mathrm{F})\) of the space \(\mathcal{L}^{\mathrm{f}}(\mathrm{E};\mathrm{F})\) of finite-rank linear maps. We denote by \(\mathcal{L}^{\mathrm{c}}(\mathrm{E};\mathrm{F})\) the space of completely continuous linear maps from E to F, and we set \(\mathcal{L}^{\mathrm{c}}(\mathrm{E}) = \mathcal{L}^{\mathrm{c}}(\mathrm{E};\mathrm{E})\) . a) The space \(\mathcal{L}^{\mathrm{c}}(\mathrm{E})\) is a closed two-sided ideal of \(\mathcal{L}(\mathrm{E})\) . Let \(u \in \mathcal{L}(\mathrm{E};\mathrm{F})\) . The map \(\ell \mapsto \ell \circ u\) is a continuous linear map from \(F'\) to \(E'\) , denoted \({}^{t}u\) . If \(u \in \mathcal{L}^{\mathrm{c}}(\mathrm{E};\mathrm{F})\) , then \({}^{t}u \in \mathcal{L}^{\mathrm{c}}(\mathrm{F}'; \mathrm{E}')\) .
\end{enumerate}

\begin{enumerate}
\def\labelenumi{\alph{enumi})}
\setcounter{enumi}{1}
\item
  Let J be a set and \(\mathrm{F} = \mathcal{C}_{0}(\mathrm{J})\) . Let \(u \in \mathcal{L}(\mathrm{E};\mathrm{F})\) . For every \(j \in J\) , the map \(\ell_{j} \colon x \mapsto u(x)(j)\) is an element of \(E'\) . We have \(\|u\| = \sup_{j \in J} \| \ell_{j} \|\) and the map \(u \mapsto (\ell_{j})_{j \in J}\) induces, by passage to subspaces, an isomorphism from \(\mathcal{L}^{\mathrm{c}}(\mathrm{E};\mathrm{F})\) to the Banach space \(\mathcal{C}_{0}(\mathrm{J};\mathrm{E}')\) .
\item
  If K is locally compact, then \(u \in \mathcal{L}(\mathrm{E}; \mathrm{F})\) is completely continuous if and only if the image under u of every bounded subset of E is relatively compact in F. (One may compare this result with prop. 11 of III, p. 16.)
\end{enumerate}

\begin{enumerate}
\def\labelenumi{\arabic{enumi})}
\setcounter{enumi}{2}
\tightlist
\item
  We denote by A the subring of K consisting of elements \(x \in \mathrm{K}\) such that \(|x| \leqslant 1\) ; it is a local ring whose maximal ideal consists of the \(x \in \mathrm{A}\) such that \(|x| < 1\) .
\end{enumerate}

Let I be a set and E the Banach space \(\mathcal{C}_{0}(\mathrm{I})\) (cf. question c) of exercise 1). For \(i \in I\) , denote by \(\varphi_{i} \in E\) the characteristic function of \(\{i\}\) . Let \(u \in \mathcal{L}^{\mathrm{c}}(\mathrm{E})\) such that \(\|u\| \leqslant 1\) .

\begin{enumerate}
\def\labelenumi{\alph{enumi})}
\item
  The subspace \(E_{0}\) of elements \(x \in E\) such that \(\|x\| \leqslant 1\) is stable under u.
\item
  For every nonzero ideal a of A, the linear map u defines, by passage to quotients, an endomorphism \(u_{a}\) of the A/a-module \(E_{a} = E_{0}/aE_{0}\) . This A/a-module is free and the image of \(u_{a}\) is contained in a finitely generated submodule.
\item
  There exists a unique formal series \(f_{u} \in A[[t]]\) such that, for every nonzero ideal a of A, the image of \(f_{u}\) in \((A/\mathfrak{a})[[t]]\) is equal to the polynomial \(\det(1 - tu_{\mathfrak{a}})\) defined in A, VIII, p. 455.
\end{enumerate}

\(d)\) Let \(v \in \mathcal{L}^{\mathrm{c}}(\mathrm{E})\) and let \(c \in \mathrm{K}^{*}\) such that \(\|cv\| \leqslant 1\) . The formal series \(f_{cv}\) is independent of the choice of \(c\) satisfying this property.

We shall denote by \(\det(1-tv)\) the formal series \(f_v\) for all \(c \in \mathbf{K}^*\) such that \(\|cv\| \leqslant 1\) .

¶ 4) We retain the notation from the previous exercise.

\begin{enumerate}
\def\labelenumi{\alph{enumi})}
\item
  Let \(v \in \mathcal{L}^c(\mathrm{E})\) . The radius of convergence of the formal series \(\det(1 - tv)\) is infinite.
\item
  Let \((v_{n})_{n\in\mathbf{N}}\) be a sequence in \(\mathcal{L}^{c}(\mathrm{E})\) converging to \(v\in\mathcal{L}^{c}(\mathrm{E})\) . The sequence of formal series \(\det(1-tv_{n})\) converges to \(\det(1-tv)\) for the topology of simple convergence of the coefficients of formal series.
\item
  Let \(v \in \mathcal{L}^{c}(E)\) . For every extension L of K and every \(x \in L\) , we denote by \(\det(1 - xv)\) the value at x of the formal power series \(\det(1 - tv)\) , which is well defined by a). We have \(\det(1 + u)\det(1 + v) = \det(1 + u + v + u \circ v)\) .
\item
  Let J be a set and \(\mathrm{F} = \mathcal{C}_{0}(\mathrm{J})\) . For all \(u \in \mathcal{L}^{\mathrm{c}}(\mathrm{E};\mathrm{F})\) and \(v \in \mathcal{L}^{\mathrm{c}}(\mathrm{F};\mathrm{E})\) , we have \(\det(1 - tu \circ v) = \det(1 - tv \circ u)\) as formal power series.
\item
  Let \(v \in \mathcal{L}^{\mathrm{c}}(\mathrm{E})\) . We denote by \(\operatorname{Tr}(v)\) the negative of the coefficient of t in \(\det(1 - tu)\) . If K is of characteristic zero, then
\end{enumerate}

\[
\det (1 - t u) = \exp \left(- \sum_ {m \geqslant 1} \frac {\operatorname{Tr} \left(u ^ {m}\right)}{m} t ^ {m}\right)
\]

as formal power series.

¶ 5) We keep the notations of the preceding exercise. Let \(v \in \mathcal{L}^{\mathrm{c}}(\mathrm{E})\) .

\begin{enumerate}
\def\labelenumi{\alph{enumi})}
\item
  The formal power series \(\det(1 - tv)\sum_{k\geqslant 0}t^{k}v^{k}\) in \(\mathcal{L}(\mathrm{E})[[t]]\) has infinite radius of convergence.
\item
  Let \(\lambda \in \mathrm{K}\) . The endomorphism \(1_{\mathrm{E}} - \lambda v\) of E is invertible if and only if \(\det(1 - \lambda v) \neq 0\) .
\item
  Let \(\lambda \in K\) such that \(\det(1 - \lambda v) = 0\) . There exist unique closed subspaces \(I_{\lambda}\) and \(N_{\lambda}\) of E, stable under \(v\) such that
\end{enumerate}

\begin{enumerate}
\def\labelenumi{(\roman{enumi})}
\item
  The space E is the topological direct sum of \(I_{\lambda}\) and \(N_{\lambda}\) ;
\item
  The endomorphism of \(I_{\lambda}\) deduced from \(1_{E} - \lambda v\) by passing to subspaces is invertible;
\item
  The endomorphism of \(N_{\lambda}\) deduced from \(1_{E} - \lambda v\) by passing to subspaces is nilpotent.
\end{enumerate}

\(d)\) Let \(h \geqslant 1\) be the order of the zero \(\lambda\) of det(1 - tv). We have \(\dim(N_{\lambda}) = h\) .

The result of this exercise is to be compared with those of No. 5 of III, p. 77, which concern Banach spaces over R or C.

In what follows, K denotes the field R or C, and all topological vector spaces are over K.

\begin{enumerate}
\def\labelenumi{\arabic{enumi})}
\setcounter{enumi}{5}
\tightlist
\item
  Let E be the space \(C^{N}\) endowed with the topology associated with the sequence of semi-norms \((p_{n})_{n\in\mathbf{N}}\) given by \(p_{n}(x)=|x_{n}|\) for every element \(x=(x_{n})_{n\in\mathbf{N}}\) of E.
\end{enumerate}

\begin{enumerate}
\def\labelenumi{\alph{enumi})}
\item
  The space E is a Fréchet space.
\item
  The set \(\mathcal{L}^{\mathrm{c}}(\mathrm{E})\) of compact endomorphisms of E is not closed in \(\mathcal{L}(\mathrm{E})\) .
\end{enumerate}

\begin{enumerate}
\def\labelenumi{\arabic{enumi})}
\setcounter{enumi}{6}
\item
  Let E and F be locally convex topological vector spaces. If F is metrizable, then the image of every compact linear map \(E \rightarrow F\) is of countable type.
\item
  Construct an example of a non-compact endomorphism \(u\) of a Hilbert space E such that \(u^2\) is compact.
\item
  For every pair \((x,y)\in]0,1]^{2}\) , there exists an integer \(n\geqslant1\) and an integer k such that \(0\leqslant k\leqslant2^{n-1}-1\) so that
\end{enumerate}

\[
2 ^ {- n} <   x \leqslant 2 ^ {- n + 1}, \quad \frac {2 k}{2 ^ {n}} <   y \leqslant \frac {2 k + 2}{2 ^ {n}}.
\]

We then define N: ]0,1]×]0,1] → C by setting

\[
\mathrm{N} (x, y) = \left\{ \begin{array}{l l} 2 & \text { if } \frac {2 k}{2 ^ {n}} <   y \leqslant \frac {2 k + 1}{2 ^ {n}} \\ 0 & \text { if } \frac {2 k + 1}{2 ^ {n}} <   y \leqslant \frac {2 k + 2}{2 ^ {n}}. \end{array} \right.
\]

\begin{enumerate}
\def\labelenumi{\alph{enumi})}
\tightlist
\item
  Show that the formula
\end{enumerate}

\[
u (f) (x) = \int_ {[ 0, 1 ]} \mathrm{N} (x, y) f (y) d y
\]

defines a continuous linear map u on \(L^{1}([0,1])\) with respect to Lebesgue measure.

\begin{enumerate}
\def\labelenumi{\alph{enumi})}
\setcounter{enumi}{1}
\item
  Show that the endomorphism \(u\) is not compact.
\item
  The endomorphism \(u^2\) is compact.
\end{enumerate}

\(d)\) Let \((x_{n})_{n\in \mathbf{N}}\) be a bounded sequence in \(\mathrm{L}^1 ([0,1])\) . The sequence \((u(x_n))_{n\in \mathbf{N}}\) admits a weakly convergent subsequence.

\begin{enumerate}
\def\labelenumi{\alph{enumi})}
\setcounter{enumi}{4}
\item
  Let \((x_{n})_{n\in \mathbf{N}}\) be a weakly convergent sequence in \(\mathrm{L}^1 ([0,1])\) . The sequence \((u(x_{n}))_{n\in \mathbf{N}}\) converges in \(\mathrm{L}^1 ([0,1])\) .
\item
  Let E be the topological vector space whose elements are the functions \(f\colon[0,1]\to\mathbf{C}\) which are bounded and measurable, equipped with the finest topology that is stronger than the topology of \(\mathrm{L}^{2}([0,1])\) and the weak topology relative to \(\mathrm{L}^{1}([0,1])\) . Then the linear map v defined by
\end{enumerate}

\[
v (f) (x) = \int_ {[ 0, 1 ]} \overline {{\mathrm{N} (y , x)}} f (y) d y
\]

is a compact endomorphism of E. Its transpose \(^{t}v\) is not compact for the strong topology on E'.

\begin{enumerate}
\def\labelenumi{\arabic{enumi})}
\setcounter{enumi}{9}
\tightlist
\item
  Let p and r be real numbers such that \(1 \leqslant p < r\) . Let u be a continuous linear map from \(\ell_{\mathrm{K}}^{r}(\mathbf{N})\) into \(\ell_{\mathrm{K}}^{p}(\mathbf{N})\) .
\end{enumerate}

We denote by \(r_{n}\) (resp. \(p_{n}\) ) the projection

\[
(x _ {i}) _ {i \in \mathbf {N}} \mapsto (x _ {0}, \ldots , x _ {n - 1}, 0, 0, \ldots)
\]

onto \(\ell_{\mathrm{K}}^{r}(\mathbf{N})\) (resp. on \(\ell_{\mathrm{K}}^{p}(\mathbf{N})\) ). We denote by \(\| \cdot \| _p\) (resp. \(\| \cdot \| _r\) ) the norm on \(\ell_{\mathrm{K}}^{p}(\mathbf{N})\) (resp. on \(\ell_{\mathrm{K}}^{r}(\mathbf{N})\) ).

\begin{enumerate}
\def\labelenumi{\alph{enumi})}
\item
  Show that for every \(n \in \mathbf{N}\) , we have \(\|p_n u(1 - r_m)\| \to 0\) as \(m \to +\infty\) where 1 denotes the identity morphism of \(\ell_{\mathrm{K}}^{r}(\mathbf{N})\) .
\item
  For \(n\) , \(m \in \mathbf{N}\) , let \(u_{n,m} = (1 - p_n)u(1 - r_m)\) . Let \(\alpha = (\alpha_n)_{n \in \mathbf{N}}\) be a sequence of strictly positive real numbers such that \(\|\alpha\|_r = 1\) and \(\alpha\) does not belong to \(\ell_{\mathrm{K}}^p(\mathbf{N})\) and let \((\varepsilon_n)\) be a sequence of strictly positive real numbers such that the series \(\sum \varepsilon_n\alpha_n\) converges.
\end{enumerate}

If \(u\) is not compact, show that there exist \(\delta > 0\) , strictly increasing sequences of integers \((n_i)_{i \in \mathbf{N}}\) and \((m_i)_{i \in \mathbf{N}}\) , and two sequences \((x_i)_{i \in \mathbf{N}}\) of elements of \(\ell_{\mathrm{K}}^r(\mathbf{N})\) and \((y_i)_{i \in \mathbf{N}}\) of elements of \(\ell_{\mathrm{K}}^p(\mathbf{N})\) with the following properties:

\begin{enumerate}
\def\labelenumi{(\roman{enumi})}
\item
  One has \(\|x_{i}\|_{r}=1\) for all \(i\in N\) ;
\item
  The sets \(S_{i} = \{n \in N \mid x_{i,n} \neq 0\}\) are pairwise disjoint and the sets \(T_{i} = \{n \in N \mid y_{i,n} \neq 0\}\) are pairwise disjoint;
\item
  One has \(\| y_i\| _p\geqslant \delta\)
\item
  One has \(\|u(x_{i})\|_{q} > \delta - 2\varepsilon_{i}\) .
\end{enumerate}

(Proceed by induction and use the formulas

\[
u = u _ {n, m} + p _ {n} u (1 - r _ {m}) + (1 - p _ {n}) u r _ {m} + p _ {n} u q _ {m}
\]

for n and m in N.)

\begin{enumerate}
\def\labelenumi{\alph{enumi})}
\setcounter{enumi}{2}
\item
  Deduce by contradiction that u is compact ("Pitt's theorem").
\item
  Show that the canonical injection of \(\ell_{\mathrm{K}}^{p}(\mathbf{N})\) into \(\ell_{\mathrm{K}}^{r}(\mathbf{N})\) is not compact.
\item
  Let p and r be real numbers such that \(1 \leqslant p \leqslant r \leqslant +\infty\) . Show that \(\ell_{\mathrm{K}}^{r}(\mathbf{N})\) is isomorphic to \(\ell_{\mathrm{K}}^{p}(\mathbf{N})\) if and only if p = r.
\end{enumerate}

\begin{enumerate}
\def\labelenumi{\arabic{enumi})}
\setcounter{enumi}{10}
\tightlist
\item
  Let E and F be Banach spaces over K. A linear map \(u \in \mathcal{L}(\mathrm{E};\mathrm{F})\) is said to be completely continuous if, for every weakly convergent sequence \((x_{n})_{n \in \mathbb{N}}\) in E, the sequence \((u(x_{n}))_{n \in \mathbb{N}}\) converges in F.
\end{enumerate}

\begin{enumerate}
\def\labelenumi{\alph{enumi})}
\item
  Every compact linear map is completely continuous.
\item
  If the dual space E' of E is countably generated, then every completely continuous linear map E → F is compact.
\item
  Let \((u_{n})_{n\in\mathbf{N}}\) be a sequence of elements of \(\ell_{\mathrm{K}}^{1}(\mathbf{N})\) . Show that \((u_{n})\) converges if and only if it converges weakly ("Schur's theorem"; cf. EVT, IV, p. 54, exercise 15).
\end{enumerate}

\(d)\) Deduce that the injection of \(\ell_{\mathrm{K}}^{1}(\mathbf{N})\) into \(\ell_{\mathrm{K}}^{2}(\mathbf{N})\) is completely continuous but is not compact.

\begin{enumerate}
\def\labelenumi{\alph{enumi})}
\setcounter{enumi}{4}
\item
  For every Banach space, every continuous linear map \(\mathrm{E} \to \ell_{\mathrm{K}}^{1}(\mathbf{N})\) is completely continuous, and every continuous linear map \(\ell_{\mathrm{K}}^{1}(\mathbf{N}) \to \mathrm{E}\) is completely continuous.
\item
  If E is reflexive or if its dual is countably generated, every continuous linear map \(\mathrm{E} \rightarrow \ell_{\mathrm{K}}^{1}(\mathbf{N})\) is compact.
\end{enumerate}

\begin{enumerate}
\def\labelenumi{\arabic{enumi})}
\setcounter{enumi}{11}
\tightlist
\item
  Let E be an infinite-dimensional reflexive Banach space. Let \(u \in \mathcal{L}^{c}(E)\) be a compact endomorphism. Suppose that u is not of finite rank. Let \(E_{\sigma}\) be the space E equipped with the weak topology.
\end{enumerate}

Show that the linear map \(^{t}u\) is a compact endomorphism of \((\mathrm{E}_{\sigma})'\) .

\begin{enumerate}
\def\labelenumi{\arabic{enumi})}
\setcounter{enumi}{12}
\tightlist
\item
  Let E be a countably generated Banach space over K. Recall (EVT, IV, p. 70, exercise 14) that a sequence \((x_{n})_{n\in \mathbf{N}}\) of elements of E is called a Schauder basis of E if, for every \(x\in \mathrm{E}\) , there exists a unique sequence of scalars \((\lambda_n(x))_{n\in \mathbf{N}}\) such that
\end{enumerate}

\[
x = \sum_ {n = 1} ^ {+ \infty} \lambda_ {n} (x) x _ {n}
\]

where the series converges in E.

\begin{enumerate}
\def\labelenumi{\alph{enumi})}
\item
  Every separable Hilbert space admits a Banach basis; every space \(\ell^{p}(\mathrm{I})\) , where I is countable, admits a Banach basis.
\item
  The Banach space E of bounded holomorphic functions on the open disk D of radius 1 in the complex plane, equipped with the norm \(\| f \| = \sup_{z \in \mathrm{D}} |f(z)|\) , has a Schauder basis.
\item
  Every Banach space E admitting a Banach basis is an approximation space. More precisely, there exists a real number \(c \geqslant 0\) such that \(1_{E}\) belongs to the closure in the space \(\mathcal{L}_{\mathrm{pc}}(\mathrm{E})\) of the set of finite-rank endomorphisms with norm \(\leqslant c\) in E.
\end{enumerate}

\begin{enumerate}
\def\labelenumi{\arabic{enumi})}
\setcounter{enumi}{13}
\tightlist
\item
  \begin{enumerate}
  \def\labelenumii{\alph{enumii})}
  \tightlist
  \item
    Let E and F be Hilbert spaces. A continuous linear map u from E into F is compact if and only if the image of u contains no closed subspace of infinite dimension.
  \end{enumerate}
\end{enumerate}

\begin{enumerate}
\def\labelenumi{\alph{enumi})}
\setcounter{enumi}{1}
\tightlist
\item
  The characterization in a) of compact endomorphisms from E into F does not hold for every pair of Banach spaces (E, F). (One may use Exercise 5, p. 121.)
\end{enumerate}

\begin{enumerate}
\def\labelenumi{\arabic{enumi})}
\setcounter{enumi}{14}
\tightlist
\item
  Let \(E_{0}\) be the complex vector space of functions \(f: N^{*} \times N^{*} \to C\) . For \(k \in N^{*}\) , let \(\alpha_{k} \in E_{0}\) be the function such that
\end{enumerate}

\[
\alpha_ {k} (n, m) = \left\{ \begin{array}{l l} m ^ {k} & \text { if } 1 \leqslant n \leqslant k - 1 \\ n ^ {k} & \text { if } n \geqslant k. \end{array} \right.
\]

Let E be the vector subspace of \(E_{0}\) formed by the functions f such that \(p_{k}(f)\) is finite for all \(k \in N^{*}\) , where \(p_{k}\) is defined by

\[
p _ {k} (f) = \sum_ {n, m} \alpha_ {k} (n, m) | f (n, m) |.
\]

\begin{enumerate}
\def\labelenumi{\alph{enumi})}
\tightlist
\item
  The space E equipped with the topology defined by the norms \((p_{k})_{k\in\mathbf{N}^{*}}\) is a Fréchet space and a Montel space. Its dual E' is identified with the subspace of \(E_{0}\) consisting of functions g such that, for at least one \(k\geqslant1\) , one has
\end{enumerate}

\[
\sup _ {m, n} \alpha_ {k} (m, n) ^ {- 1} | g (m, n) | <   + \infty
\]

(cf. EVT, IV, p. 63, exercise 8).

\begin{enumerate}
\def\labelenumi{\alph{enumi})}
\setcounter{enumi}{1}
\tightlist
\item
  Consider the linear map \(u\) from E into \(\ell^1 (\mathbf{N}^*)\) defined by
\end{enumerate}

\[
u (f) = \left(\sum_ {n} f (n, m)\right) _ {m \in \mathbf {N} ^ {*}}.
\]

The linear map u is continuous.

\begin{enumerate}
\def\labelenumi{\alph{enumi})}
\setcounter{enumi}{2}
\item
  Identify \(\ell^{\infty}(\mathbf{N}^{*})\) with the dual of \(\ell^{1}(\mathbf{N}^{*})\) . The transpose of \(u\) is the map \(v\colon \ell^{\infty}(\mathbf{N}^{*})\to \mathrm{E}^{\prime}\subset \mathrm{E}_{0}\) determined by \((x_{n})\mapsto f\) where \(f(n,m) = x_{m}\) for all \((n,m)\in \mathbf{N}^{*}\times \mathbf{N}^{*}\) .
\item
  The map u induces an isomorphism \(\mathrm{E}/\mathrm{Ker}(u)\to\ell^{1}(\mathbf{N}^{*})\) (cf. loc. cit.).
\item
  The map u is not compact.
\item
  The transpose v of u is compact for the topology of convergence on bounded subsets of \(\ell^{\infty}(\mathbf{N}^{*})\) .
\item
  The map v is an example of a compact linear map between separated locally convex spaces whose transpose is not compact.
\end{enumerate}

\begin{enumerate}
\def\labelenumi{\arabic{enumi})}
\setcounter{enumi}{15}
\item
  Let E be a locally convex space and u a compact endomorphism of E. Denote by \(E_{b}^{\prime}\) the dual space of E equipped with the topology of bounded convergence. The linear map \(^{t}u \circ ^{t}u\) is a compact endomorphism of \(E_{b}^{\prime}\) .
\item
  Let E be a Banach space of countable type. Then \(\mathcal{L}^{\mathrm{c}}(\mathrm{E})\) , equipped with the topology induced by the norm of \(\mathcal{L}(\mathrm{E})\) is of countable type if, and only if, the dual space \(\mathrm{E}'\) is of countable type.
\item
  Let E, F, G be Banach spaces and let u be an element of \(\mathcal{L}^{c}(\mathrm{E};\mathrm{F})\) .
\end{enumerate}

\begin{enumerate}
\def\labelenumi{\alph{enumi})}
\tightlist
\item
  Let v be an element of \(\mathcal{L}(\mathrm{E};\mathrm{G})\) . Suppose that, for every \(\varepsilon > 0\) , there exists C \textgreater{} 0 such that
\end{enumerate}

\[
\| v (x) \| _ {\mathrm{G}} \leqslant \varepsilon \| x \| _ {\mathrm{E}} + \mathrm{C} \| u (x) \| _ {\mathrm{F}}
\]

for all \(x\in \mathrm{E}\) . Then \(v\in \mathcal{L}^{\mathrm{c}}(\mathrm{E};\mathrm{G})\) .

\begin{enumerate}
\def\labelenumi{\alph{enumi})}
\setcounter{enumi}{1}
\tightlist
\item
  Let w be an injective element of \(\mathcal{L}(\mathrm{F};\mathrm{G})\) . Then, for every \(\varepsilon > 0\) , there exists C \textgreater{} 0 such that
\end{enumerate}

\[
\| u (x) \| _ {\mathrm{F}} \leqslant \varepsilon \| x \| _ {\mathrm{E}} + \mathrm{C} \| (w \circ u) (x) \| _ {\mathrm{G}}
\]

for all \(x \in E\) .

\begin{enumerate}
\def\labelenumi{\arabic{enumi})}
\setcounter{enumi}{18}
\tightlist
\item
  Let X be a locally compact, metrizable, and countable at infinity topological space, and \(\mu\) a positive measure on X. We denote by A the set of atoms of \(\mu\) . Let \(f \in \mathcal{L}^{\infty}(X, \mu)\) and \(p \in [1, +\infty]\) .
\end{enumerate}

Show that multiplication by \(f\) defines a compact endomorphism of \(\mathrm{L}^{p}(\mathrm{X},\mu)\) if and only if the following two conditions are satisfied:

\begin{enumerate}
\def\labelenumi{(\roman{enumi})}
\item
  The set \(\{x\in X-A\mid f(x)\neq0\}\) is \(\mu\) -negligible;
\item
  For every \(\varepsilon > 0\) there exists a finite subset F of A such that \(|f(a)| \leqslant \varepsilon\) for all \(a \in \mathbf{A} - \mathbf{F}\) .
\end{enumerate}

(Consider multiplication by \(f\) on the space \(\mathrm{L}^{p}(\mathrm{X}_{\varepsilon},\mu)\) , where \(\mathrm{X}_{\varepsilon}\) is the set of \(x\in\mathrm{X}\) such that \(|f(x)|\geqslant\varepsilon\) .)

\begin{enumerate}
\def\labelenumi{\arabic{enumi})}
\setcounter{enumi}{19}
\tightlist
\item
  Let X, Y be compact metric spaces, T a compact operator from \(\mathcal{C}(\mathrm{Y})\) into \(\mathcal{C}(\mathrm{X})\) .
\end{enumerate}

Prove that there exists a positive measure \(\mu\) on \(X\) and a measurable map \(N\) from \(X \times Y\) into \(\mathbf{C}\) such that

\begin{enumerate}
\def\labelenumi{\alph{enumi})}
\item
  For every \(y \in Y\) , the map \(\mathrm{N}_{y} \colon x \mapsto \mathrm{N}(x, y)\) from \(X\) into \(\mathbf{C}\) is \(\mu\) -integrable;
\item
  The map \(y \mapsto N_{y}\) from Y into \(\mathcal{L}^{1}(X, \mu)\) is continuous;
\item
  For every \(f \in \mathcal{C}(\mathrm{X})\) , and every \(y \in \mathrm{Y}\) , we have
\end{enumerate}

\[
\mathrm{T} f (y) = \int_ {\mathrm{X}} \mathrm{N} (x, y) f (x) d \mu (x).
\]

\begin{enumerate}
\def\labelenumi{\arabic{enumi})}
\setcounter{enumi}{20}
\tightlist
\item
  We denote by D the open disk with center 0 and radius 1 in C, and by H \(^{2}\) (D) the Hardy space on D (EVT, V, p. 7, example 5). We write T = R/2πZ, which we equip with the normalized Haar measure. Let j be the isometry from H \(^{2}\) (D) into L \(^{2}\) (T) defined by
\end{enumerate}

\[
j (f) (\theta) = \sum_ {n \geqslant 0} \frac {f ^ {(n)} (0)}{n !} e ^ {i n \theta}.
\]

The image of H \(^{2}\) (D) by \(j\) will be denoted H \(^{2}\) (T).

Let \(\widetilde{\mathrm{P}}\) be the orthogonal projector of \(\mathrm{L}^2 (\mathbf{T})\) with image \(\mathrm{H}^2 (\mathbf{T})\), and let P be the linear map \(f\mapsto \widetilde{\mathrm{P}} (f)\) from \(\mathrm{L}^2 (\mathbf{T})\) into \(\mathrm{H}^2 (\mathbf{T})\) . If \(f\) is an element of \(\mathrm{L}^{\infty}(\mathbf{T})\), let \(\mathrm{M}_f\) be the endomorphism \(g\mapsto fg\) of \(\mathrm{L}^2 (\mathbf{T})\) .

Show that if \(f\) is a continuous function, then the commutator \(\widetilde{\mathrm{P}}\mathrm{M}_f - \mathrm{M}_f\widetilde{\mathrm{P}}\) is a compact linear map. (Consider the particular case where \(f\) is a trigonometric polynomial.)

\begin{enumerate}
\def\labelenumi{\arabic{enumi})}
\setcounter{enumi}{21}
\tightlist
\item
  We keep the notations from the previous exercise. Let R be the continuous linear map from \(\mathrm{H}^{2}(\mathbf{T})\) into \(\mathrm{L}^{2}(\mathbf{T})\) defined by \(\mathrm{R}f(\theta)=f(-\theta)\) . For \(g\in\mathrm{L}^{\infty}(\mathbf{T})\) we set \(\Gamma_{g}=\mathrm{PM}_{g}\mathrm{R}\) .
\end{enumerate}

\begin{enumerate}
\def\labelenumi{\alph{enumi})}
\item
  Prove that the linear map \(\Gamma\colon g\mapsto\Gamma_{g}\) from \(\mathrm{L}^{\infty}(\mathbf{T})\) into \(\mathcal{L}(\mathrm{H}^{2}(\mathbf{T}))\) is continuous, with norm \(\leqslant 1\), and its kernel F is equal to \(\mathrm{H}^{2}(\mathbf{T})^{\circ}\cap\mathrm{L}^{\infty}(\mathbf{T})\) .
\item
  If \(g\) is a continuous function, then \(\Gamma_g\) is a compact endomorphism of \(\mathrm{H}^2 (\mathbf{T})\) .
\item
  Let \(g \in \mathrm{L}^{\infty}(\mathbf{T})\) . We denote by \(\alpha_{g}\) the linear form on \(\mathrm{H}^{2}(\mathbf{T})\) defined by
\end{enumerate}

\[
\alpha_ {g} (f) = \int_ {\mathbf {T}} \Gamma_ {g} (f) (\theta) d \theta .
\]

Let \(\mathcal{A}\) be the image under \(j\) of the subspace of \(\mathrm{H}^2 (\mathrm{D})\) consisting of holomorphic functions that extend analytically to a neighborhood of \(\overline{\mathrm{D}}\) . Prove that if \(f_{1}\) and \(f_{2}\) belong to \(\mathcal{A}\) , then

\[
\alpha_ {g} (f _ {1} f _ {2}) = \int_ {\mathbf {T}} \Gamma_ {g} (f _ {1}) (\theta) f _ {2} (- \theta) d \theta .
\]

\(d)\) Let \(\varphi\) be a holomorphic function in a neighborhood of \(\overline{\mathrm{D}}\) not vanishing on \(\partial\mathrm{D}\) and let \(f = j(\varphi|\mathrm{D})\) . Then there exist \(f_{1}\) and \(f_{2}\) in \(\mathscr{A}\) such that \(f = f_{1}f_{2}\) and

\[
\| f \| _ {\mathrm{L} ^ {1} (\mathbf {T})} = \| f _ {1} \| _ {\mathrm{L} ^ {2} (\mathbf {T})} \| f _ {2} \| _ {\mathrm{L} ^ {2} (\mathbf {T})}.
\]

(Let \(z_{1},\ldots,z_{n}\) be the zeros of \(\varphi\) in D; then there exists a function \(\psi\) holomorphic in a neighborhood of D such that

\[
\varphi (z) = \prod_ {k = 1} ^ {n} \frac {z - z _ {k}}{1 - \overline {{z}} _ {k} z} \psi (z) ^ {2}
\]

for all z.)

\begin{enumerate}
\def\labelenumi{\alph{enumi})}
\setcounter{enumi}{4}
\tightlist
\item
  Prove that, for all \(f \in \mathrm{H}^{2}(\mathbf{T})\) , we have the inequality
\end{enumerate}

\[
\left| \alpha_ {g} (f) \right| \leqslant \left\| \Gamma_ {g} \right\| _ {\mathscr {L} \left(\mathrm{H} ^ {2} (\mathbf {T})\right)} \| f \| _ {\mathrm{L} ^ {1} (\mathbf {T})}.
\]

(If \(\varphi \in \mathrm{H}^2 (\mathrm{D})\) , there exists a sequence \((r_n)_{n\in \mathbf{N}}\) of ]0,1[ such that, for every \(n\in \mathbf{N}\) , the function \(\varphi_{n}\) defined by \(\varphi_{n}(z) = \varphi (r_{n}z)\) has no zeros on \(\partial \mathrm{D}\) , and \(\lim \varphi_n = \varphi\) in \(\mathrm{H}^2 (\mathrm{D}).\)

Deduce that there exists \(h \in \mathrm{L}^{\infty}(\mathbf{T})\) such that

\[
\alpha_ {g} (f) = \int_ {\mathbf {T}} f (\theta) h (\theta) d \theta
\]

for all \(f \in \mathrm{H}^2(\mathbf{T})\) , and \(\|h\|_{\mathrm{L}^\infty(\mathbf{T})} \leqslant \| \Gamma_g \|_{\mathcal{L}(\mathrm{H}^2(\mathbf{T}))}\) .

\(f)\) Prove that there exists \(k \in \mathrm{L}^{\infty}(\mathbf{T})\) such that \(\Gamma_g = \Gamma_k\) and \(\| k \|_{\mathrm{L}^{\infty}(\mathbf{T})} = \| \Gamma_g \|_{\mathcal{L}(\mathrm{H}^2(\mathbf{T}))}\) .

The map \(\widehat{\Gamma}\colon\mathrm{L}^{\infty}(\mathbf{T})/\mathrm{F}\to\mathcal{L}(\mathrm{H}^{2}(\mathbf{T}))\) is therefore an isometry ("Nehari's theorem").

\(g)\) Let \(g\) be a continuous function on \(\mathbf{T}\) . Then we have \(d(g, \mathrm{F}) = d(g, \mathrm{F} \cap \mathcal{C}(\mathbf{T}))\) . (For \(\mathrm{R} > 1\) , one defines the linear map \(\mathrm{P}_{\mathrm{R}}\) from \(\mathrm{L}^{\infty}(\mathbf{T})\) into \(\mathcal{C}(\mathbf{T})\) by

\[
\mathrm{P} _ {\mathrm{R}} f (\theta) = \int_ {\mathbf {T}} \frac {\mathrm{R} ^ {2} - 1}{| \mathrm{R} e ^ {i (\theta - \varphi)} - 1 | ^ {2}} f (\varphi) d \varphi .
\]

Then \(\| \mathrm{P}_{\mathrm{R}}f\|_{\mathcal{C}(\mathbf{T})}\leqslant \| f\|_{\mathrm{L}^{\infty}(\mathbf{T})}\) and, if \(f\in \mathcal{C}(\mathbf{T})\) , we have \(\mathrm{P_R}f\to f\) in \(\mathcal{C}(\mathbf{T})\) as \(\mathrm{R}\rightarrow 1.\)

Deduce that \(\Gamma(\mathcal{C}(\mathbf{T}))\) is a closed subspace of \(\mathcal{L}(\mathrm{H}^{2}(\mathbf{T}))\) ("Sarason's theorem").

\begin{enumerate}
\def\labelenumi{\alph{enumi})}
\setcounter{enumi}{7}
\tightlist
\item
  Let \(g \in \mathrm{L}^{\infty}(\mathbf{T})\) such that \(\Gamma_{g}\) is a compact endomorphism. For \(m \in Z\) we set \(e_{m}(\theta) = e^{im\theta}\) . Prove that \(\|\Gamma_{ge_{-m}}\|_{\mathcal{L}(\mathrm{H}^{2}(\mathbf{T}))}\) tends to zero as m tends to \(+\infty\) , and then that there exists a sequence \((h_{m})_{m \in \mathbf{N}}\) of F such that the sequence \((e_{m}h_{m})_{m \in \mathbf{N}}\) converges to g in \(\mathrm{L}^{\infty}(\mathbf{T})\) . Deduce that there exists a continuous function in the class of g modulo F ("Hartman's theorem").
\end{enumerate}

\begin{enumerate}
\def\labelenumi{\arabic{enumi})}
\setcounter{enumi}{22}
\item
  Let B be an infinite-dimensional Banach space. We denote by E the strong dual of B and by F the dual of B equipped with the topology \(\sigma(\mathrm{B}^{\prime},\mathrm{B})\) . Prove that the identity map of the dual of B is compact from E into F, but that its transpose is not compact from \(\mathrm{F}_{b}^{\prime}\) into \(\mathrm{E}_{b}^{\prime}\) .
\item
  Let E be a real or complex Banach space and \(u \in \mathcal{L}(\mathrm{E})\) . Assume that the sequence \((u^{n})_{n \in \mathbf{N}}\) is bounded and that u is compact from E into E equipped with the weak topology \(\sigma(\mathrm{E}, \mathrm{E}')\) . For every integer \(N \geqslant 1\) we set
\end{enumerate}

\[
v _ {\mathrm{N}} = \frac {1}{\mathrm{N}} (1 _ {\mathrm{E}} + u + \dots + u ^ {\mathrm{N-1}}) \in \mathcal {L} (\mathrm{E}).
\]

\begin{enumerate}
\def\labelenumi{\alph{enumi})}
\item
  Let F be the closure of the image of \(1_{\mathrm{E}} - u\) . For every \(x \in \mathrm{F}\) , we have \(v_{\mathrm{N}}(x) \to 0\) in E when \(\mathrm{N} \to +\infty\) .
\item
  Let \(x \in E\) . The sequence \((v_{\mathrm{N}}(x))\) converges to an element \(y \in E\) such that \(u(y) = y\) . (There exists a subsequence of \((v_{\mathrm{N}}(x))\) that converges weakly to \(y \in E\) ; show that \(x - y \in F\) using the Hahn--Banach theorem.)
\item
  The map \(w \colon \mathrm{E} \to \mathrm{E}\) such that \(w(x)\) is the element \(y\) given by the preceding question is a continuous projector whose image is the eigenspace of \(u\) corresponding to the eigenvalue 1.
\end{enumerate}

\(d)\) Let \(z \in \mathbf{C}\) such that \(|z| = 1\) . The sequence in \(\mathcal{L}(\mathrm{E})\) defined by

\[
v _ {z, \mathrm{N}} = \frac {1}{\mathrm{N}} (1 _ {\mathrm{E}} + z ^ {- 1} u \dots + z ^ {- (\mathrm{N-1})} u ^ {\mathrm{N-1}})
\]

converges in \(\mathcal{L}(\mathrm{E})\) equipped with the topology of pointwise convergence. Let \(w_{z}\) be its limit. Then \(w_{z}\) is a projector, one has \(w_{z}w = ww_{z} = zw_{z}\) and \(w_{z_1}w_{z_2} = 0\) if \(z_{1} \neq z_{2}\) .

\begin{enumerate}
\def\labelenumi{\arabic{enumi})}
\setcounter{enumi}{24}
\tightlist
\item
  \begin{enumerate}
  \def\labelenumii{\alph{enumii})}
  \tightlist
  \item
    Let \(N \geqslant 1\) be an integer. We denote by \(G_{N}\) (resp. \(H_{N}\) ) the set \(\{-1,1\}^{N}\) (resp. the set \(\{-1,2\}^{N}\) ). We denote by \(\mu_{N}\) (resp. \(\nu_{N}\) ) the positive measure of total mass 1 on \(G_{N}\) (resp. on \(H_{N}\) ) such that, for every \((t_{i})_{1 \leqslant i \leqslant N}\) in \(G_{N}\) (resp. in \(H_{N}\) ), one has \(\mu_{\mathrm{N}}((t_{i})_{1 \leqslant i \leqslant N}) = \frac{1}{2^{\mathrm{N}}}\) , resp.
  \end{enumerate}
\end{enumerate}

\[
\nu_ {\mathrm{N}} ((t _ {i}) _ {1 \leqslant i \leqslant \mathrm{N}}) = \left(\frac {1}{3}\right) ^ {j} \left(\frac {2}{3}\right) ^ {\mathrm{N} - j},
\]

where j is the number of indices i such that \(t_{i}=2\) .

There exists a real number \(c > 0\) such that, for every family \((a_{i})_{1\leqslant i\leqslant N}\) of complex numbers, we have

\[
\mu_ {\mathrm{N}} \Big (\Big \{(t _ {i}) _ {1 \leqslant i \leqslant \mathrm{N}} \in \mathrm{G} _ {\mathrm{N}} | \Big | \sum_ {i = 1} ^ {\mathrm{N}} a _ {i} t _ {i} \Big | > c (\log \mathrm{N}) ^ {1 / 2} \Big (\sum_ {i = 1} ^ {\mathrm{N}} | a _ {i} | ^ {2} \Big) ^ {1 / 2} \Big \} \Big) <   \frac {c}{\mathrm{N} ^ {3}},
\]

and

\[
\nu_ {\mathrm{N}} \Big (\Big \{(t _ {i}) _ {1 \leqslant i \leqslant \mathrm{N}} \in \mathrm{H} _ {\mathrm{N}} | \left| \sum_ {i = 1} ^ {\mathrm{N}} a _ {i} t _ {i} \right| > c (\log \mathrm{N}) ^ {1 / 2} \Big (\sum_ {i = 1} ^ {\mathrm{N}} | a _ {i} | ^ {2} \Big) ^ {1 / 2} \Big \} \Big) <   \frac {c}{\mathrm{N} ^ {3}}.
\]

(We may assume that \(a_{i} \in \mathbf{R}\) for all \(i\) and that \(\sum |a_{i}|^{2} = 1\) . For the first inequality, prove and use the upper bound.

\[
\int_ {G _ {N}} \exp \left(\lambda \left| \sum_ {i = 1} ^ {N} a _ {i} t _ {i} \right|\right) d \mu_ {N} \leqslant 2 \exp (\lambda^ {2})
\]

for every real \(\lambda > 0\) ; reason in a similar manner for the second upper bound using the inequality \(\frac{1}{3}(e^{2x} + 2e^{-x}) \leqslant e^{2x^2}\) for all \(x \in \mathbf{R}\) .)

\begin{enumerate}
\def\labelenumi{\alph{enumi})}
\setcounter{enumi}{1}
\tightlist
\item
  For \(k \in N\) , let \(\mathrm{A}_{k} = \mathbf{Z}/(3 \cdot 2^{k})\mathbf{Z}\) . There exists a real number \(c' > 0\) , not depending on k, and a function \(f_{k} \colon \widehat{\mathrm{A}}_{k} \to \{-1, 2\}\) such that
\end{enumerate}

\[
\sum_ {\chi \in \widehat {\mathrm{A}} _ {k}} f _ {k} (\chi) = 0
\]

and, for all \(x \in A_{k}\) , one has

\[
\Big | \sum_ {\chi \in \widehat {\mathrm{A}} _ {k}} \chi (x) f _ {k} (\chi) \Big | \leqslant c ^ {\prime} (k + 1) ^ {1 / 2} 2 ^ {k / 2}.
\]

(Using \(a\) ), determine the existence of a function \(f_{k}\) satisfying the second condition, then modify \(f_{k}\) to obtain the first.)

We denote by \(B_{k}\) (resp. \(C_{k}\) ) the set of characters \(\chi\) of \(A_{k}\) such that \(f_{k}(\chi)=2\) (resp. \(f_{k}(\chi)=-1\) ). We have \(\mathrm{Card(B_{k})}=2^{k}\) and \(\mathrm{Card(C_{k})}=2^{k+1}\) .

For every \(k \geqslant 1\) , we fix a bijection \(\iota_k \colon B_k \to C_{k-1}\) .

We denote by A the disjoint union of the \(A_{k}\) for \(k \in N\) ; we equip A with the discrete topology, and we denote by \(\mathrm{E} = \mathcal{C}_{b}(\mathrm{A})\) .

We fix a family \(\varepsilon = (\varepsilon_k)_{k\geqslant \mathbf{N}}\) of functions \(\varepsilon_{k}\colon \mathrm{B}_{k}\to \{-1,1\}\) . We then define functions \(e_{k,\chi}\in \mathrm{E}\) , for \(k\in \mathbf{N}\) and \(\chi \in \mathrm{B}_k\) by

\[
e _ {k, \chi} (x) = \left\{ \begin{array}{l l} \iota_ {k} (\chi) (x) & \text { if } k \geqslant 1 \text {   and   } x \in \mathrm{A} _ {k - 1}, \\ \varepsilon_ {k} (\chi) \chi (x) & \text { if } x \in \mathrm{A} _ {k} \\ 0 & \text { otherwise. } \end{array} \right.
\]

We denote by \(E_{\varepsilon}\) the closed subspace of E generated by the functions \((e_{k,\chi})_{k\in\mathbf{N},\chi\in\mathbf{B}_{k}}\) . This is a complex Banach space of countable type. The rest of the exercise aims to show that, for a suitable choice of \(\varepsilon\) , the space \(\mathrm{E}_{\varepsilon}\) does not have the approximation property; it admits no Banach basis.

\begin{enumerate}
\def\labelenumi{\alph{enumi})}
\setcounter{enumi}{2}
\tightlist
\item
  For \(k \in N\) and \(\chi \in B_{k}\) there exists a continuous linear form \(\lambda_{k,\chi} \in E_{\varepsilon}'\) such that, for every \(f \in E_{\varepsilon}\) ,
\end{enumerate}

\[
\lambda_ {k, \chi} (f) = \frac {1}{\operatorname{Card} (\mathrm{A} _ {k})} \sum_ {x \in \mathrm{A} _ {k}} \varepsilon_ {k} (\chi) \chi (x ^ {- 1}) f (x).
\]

For \(k \geqslant 1\) and \(f \in \mathrm{E}_{\varepsilon}\) , we have

\[
\lambda_ {k, \chi} (f) = \frac {2}{\operatorname{Card} (\mathrm{A} _ {k})} \sum_ {x \in \mathrm{A} _ {k}} \iota_ {k} (\chi) (x ^ {- 1}) f (x).
\]

\begin{enumerate}
\def\labelenumi{\alph{enumi})}
\setcounter{enumi}{3}
\tightlist
\item
  For every \(k \in N\) , let \(\beta_{k}\) be the continuous linear form on \(\mathcal{L}(E_{\varepsilon})\) defined by
\end{enumerate}

\[
\beta_ {k} (u) = \frac {1}{2 ^ {k}} \sum_ {\chi \in \mathrm{B} _ {k}} \lambda_ {k, \chi} (u (e _ {k, \chi})).
\]

One has \(\beta_{k}(1_{\mathrm{E}_{\varepsilon}}) = 1\)

\begin{enumerate}
\def\labelenumi{\alph{enumi})}
\setcounter{enumi}{4}
\tightlist
\item
  Let \(u \in \mathcal{L}(\mathrm{E}_{\varepsilon})\) and \(k \in \mathbf{N}\) . We have
\end{enumerate}

\[
\beta_ {k + 1} (u) - \beta_ {k} (u) = \frac {1}{\operatorname{Card} \left(\mathrm{A} _ {k}\right)} \sum_ {x \in \mathrm{A} _ {k}} u \left(\varphi_ {k, x}\right) (x)
\]

where, for \(x \in \mathrm{A}_k\) , we set

\[
\varphi_ {k, x} = \frac {1}{2 ^ {k + 1}} \sum_ {\chi \in B _ {k + 1}} \iota_ {k + 1} (\chi) (x ^ {- 1}) e _ {k + 1, \chi} - \frac {1}{2 ^ {k}} \sum_ {\chi \in B _ {k}} \varepsilon_ {k} (\chi) \chi (x ^ {- 1}) e _ {k, \chi}.
\]

One has \(|\varphi_{k,x}(y)| \leqslant c'(k + 1)^{1/2}2^{-k/2}\) for \((x,y) \in \mathrm{A}_k^2\) .

\begin{enumerate}
\def\labelenumi{\alph{enumi})}
\setcounter{enumi}{5}
\item
  There exists a real number \(c'' \geqslant c'\) such that, for every \(k \geqslant 1\) , there exists a function \(\varepsilon_k \colon B_k \to \{-1,1\}\) satisfying \(|\varphi_{k,x}(y)| \leqslant c''3(k+1)^{1/2}2^{-k/2}\) for \((x,y) \in A_k \times A_{k-1}\) . (Apply the first inequality of a).) We now assume that \(\varepsilon = (\varepsilon_k)\) satisfies these properties.
\item
  For every \(k \in \mathbf{N}\) and \(x \in \mathrm{A}_k\) , we have \(\| \varphi_{k,x} \| \leqslant c''(k + 1)^{1/2} 2^{-k/2}\) .
\item
  Let K be the subset of \(E_{\varepsilon}\) whose elements are \(e_{0,\chi_{0}}\) , where \(\chi_{0}\) is the unique element of \(B_{0}\) , and the functions \((k+1)^{2}\varphi_{k,x}\) for \(k\in N\) and \(x\in A_{k}\) . The set K is relatively compact in \(E_{\varepsilon}\) , and we have
\end{enumerate}

\[
| \beta_ {k + 1} (u) - \beta_ {k} (u) | \leqslant \frac {1}{(k + 1) ^ {2}} \sup _ {x \in K} \| u (x) \|
\]

for all \(k \in \mathbf{N}\) and every \(u \in \mathcal{L}(\mathrm{E}_{\varepsilon})\) .

\begin{enumerate}
\def\labelenumi{\roman{enumi})}
\tightlist
\item
  The sequence \((\beta_{k})_{k\in \mathbf{N}}\) converges in the dual of the space \(\mathcal{L}(\mathrm{E}_{\varepsilon})\) equipped with the topology of pointwise convergence. Denote by \(\beta\) its limit. For every \(u\in \mathcal{L}(\mathrm{E}_{\varepsilon})\) , we have \(|\beta (u)|\leqslant 3\sup_{x\in \mathrm{K}}\| u(x)\|\) .
\end{enumerate}

\begin{enumerate}
\def\labelenumi{\alph{enumi})}
\setcounter{enumi}{9}
\tightlist
\item
  One has \(\beta(1_{\mathrm{E}_{\varepsilon}})=1\) and \(\beta(u)=0\) if \(u\) has finite rank.
\end{enumerate}

\(k)\) The space \(\mathrm{E}_{\varepsilon}\) does not have the approximation property.

\begin{enumerate}
\def\labelenumi{\alph{enumi})}
\setcounter{enumi}{11}
\tightlist
\item
  Let \(p \in]2, +\infty[\) . We equip A with the counting measure and denote by \(E_{\varepsilon}^{p}\) the closed subspace of \(L^{p}(A)\) generated by the \((e_{k,\chi})_{k \in \mathbf{N}, \chi \in \mathrm{B}_{k}}\) . The space \(E_{\varepsilon}^{p}\) does not have the approximation property. (Verify that the cardinality of the support of \(\varphi_{k,x}\) is \(\leqslant 21 \cdot 2^{k}\) and deduce that \(\|\varphi_{k,x}\|_{p} = \mathrm{O}((k+1)^{1/2}2^{-k(p-2)/2p})\) for \(k \in N\) and \(x \in A_{k}\) .)
\end{enumerate}

\begin{enumerate}
\def\labelenumi{\arabic{enumi})}
\setcounter{enumi}{25}
\tightlist
\item
  Let E and F be Banach spaces, and let u be a continuous linear map from E into F.
\end{enumerate}

\begin{enumerate}
\def\labelenumi{\alph{enumi})}
\tightlist
\item
  Prove that the following conditions are equivalent:
\end{enumerate}

\begin{enumerate}
\def\labelenumi{(\roman{enumi})}
\item
  The map u is compact;
\item
  There exists a sequence \((x_n')_n\) tending to 0 in the strong dual \(\mathrm{E}'\) of E satisfying \(\| u(x)\| \leqslant \sup |\langle x,x_n' \rangle|\) for all \(x \in \mathrm{E}\) ;
\item
  There exists a closed subspace L of the Banach space \(c_{0}\) and compact linear maps \(f: E \to L\) and \(g: L \to F\) such that \(u = g \circ f\) (If u is compact, apply cor. 1 of EVT, IV, p. 24 to the image under \(^{t}u\) of the polar of the unit ball of F. Under hypothesis (ii), let \((\lambda_{n})\) be an element of \(c_{0}\) such that the sequence \((\lambda_{n}^{-1}x_{n}^{\prime})\) tends to 0; define f by \(f(x) = (\langle x, \lambda_{n}^{-1}x_{n}^{\prime} \rangle)_{n}\) and take for L the closure of the image of f.)
\end{enumerate}

\begin{enumerate}
\def\labelenumi{\alph{enumi})}
\setcounter{enumi}{1}
\tightlist
\item
  If \(u\) is compact, we have \(\|u\| = \inf(\sup_n \|x_n'\|)\) where the infimum is taken over the set of sequences \((x_n')\) satisfying (ii), and \(\|u\| = \inf(\|f\|\|g\|)\) , the infimum being taken over the set of triples (L, \(f,g\) ) satisfying the conditions of (iii).
\end{enumerate}

\begin{enumerate}
\def\labelenumi{\arabic{enumi})}
\setcounter{enumi}{26}
\tightlist
\item
  Let E and F be locally convex spaces. A mapping \(u \in \mathcal{L}(\mathrm{E};\mathrm{F})\) is said to be weakly compact if it is a compact mapping from E into F equipped with the topology \(\sigma(\mathrm{F},\mathrm{F}')\) .
\end{enumerate}

\begin{enumerate}
\def\labelenumi{\alph{enumi})}
\item
  The weakly compact mappings form a subspace \(\mathcal{L}^{\mathrm{fc}}(\mathrm{E};\mathrm{F})\) of \(\mathcal{L}(\mathrm{E};\mathrm{F})\) . Let \(\mathrm{E}_1\) , \(\mathrm{F}_1\) be locally convex spaces, \(f\in \mathcal{L}(\mathrm{E}_1;\mathrm{E})\) , \(g\in \mathcal{L}(\mathrm{F};\mathrm{F}_1)\) and \(u\in \mathcal{L}(\mathrm{E};\mathrm{F})\) ; if \(u\) is weakly compact, the same is true of \(g\circ u\circ f\) .
\item
  One has \(\mathcal{L}^{c}(\mathrm{E};\mathrm{F})\subset\mathcal{L}^{\mathrm{fc}}(\mathrm{E};\mathrm{F})\) . If E is an infinite-dimensional reflexive Banach space, the identity map of E is weakly compact but is not compact.
\end{enumerate}

\begin{enumerate}
\def\labelenumi{\arabic{enumi})}
\setcounter{enumi}{27}
\tightlist
\item
  Let E and F be locally convex spaces and \(u \in \mathcal{L}(\mathrm{E};\mathrm{F})\) .
\end{enumerate}

\begin{enumerate}
\def\labelenumi{\alph{enumi})}
\tightlist
\item
  Show that the following properties are equivalent:
\end{enumerate}

\begin{enumerate}
\def\labelenumi{(\roman{enumi})}
\item
  The image under \(u\) of every bounded subset of E is relatively compact for the weak topology;
\item
  The map \(^{tt}u\) maps F'' into E.
\end{enumerate}

\begin{enumerate}
\def\labelenumi{\alph{enumi})}
\setcounter{enumi}{1}
\tightlist
\item
  If F is quasi-complete, prove that these properties are equivalent to:
\end{enumerate}

\begin{enumerate}
\def\labelenumi{(\roman{enumi})}
\setcounter{enumi}{2}
\tightlist
\item
  The map \(^{t}u\) maps every equicontinuous subset of F' to a relatively compact subset of E' for \(\sigma(\mathrm{E}',\mathrm{E}'')\) .
\end{enumerate}

\begin{enumerate}
\def\labelenumi{\alph{enumi})}
\setcounter{enumi}{2}
\tightlist
\item
  If E and F are Banach spaces, (i) to (iii) are also equivalent to
\end{enumerate}

(i') The map u is weakly compact;

(iii') The map \(^{t}u\colon F_{b}^{\prime}\to E_{b}^{\prime}\) is weakly compact.

\begin{enumerate}
\def\labelenumi{\arabic{enumi})}
\setcounter{enumi}{28}
\tightlist
\item
  \begin{enumerate}
  \def\labelenumii{\alph{enumii})}
  \tightlist
  \item
    Prove that every weakly compact subset of the Banach space \(\ell^{1}(\mathbf{N})\) is compact (cf. EVT, IV, p. 69, exerc. 11). In particular, every weakly compact linear map from a locally convex space E into \(\ell^{1}(\mathbf{N})\) is compact.
  \end{enumerate}
\end{enumerate}

\begin{enumerate}
\def\labelenumi{\alph{enumi})}
\setcounter{enumi}{1}
\tightlist
\item
  Let E be a Banach space. Deduce from a) that every weakly compact linear map from \(c_{0}\) into E is compact.
\end{enumerate}

\begin{enumerate}
\def\labelenumi{\arabic{enumi})}
\setcounter{enumi}{29}
\tightlist
\item
  \begin{enumerate}
  \def\labelenumii{\alph{enumii})}
  \tightlist
  \item
    Let E and F be Banach spaces and \(u \in \mathcal{L}(\mathrm{E};\mathrm{F})\) . Show that the following properties are equivalent:
  \end{enumerate}
\end{enumerate}

\begin{enumerate}
\def\labelenumi{(\roman{enumi})}
\item
  The image under \(u\) of every weakly compact subset of E is (strongly) compact;
\item
  The image under \(u\) of every weakly relatively compact subset of E is relatively compact in F;
\item
  The image under \(u\) of every weakly convergent sequence in E converges strongly in F.
\end{enumerate}

(Use the Šmulian theorem, EVT, IV, p. 36, th. 2.)

One says that u is semi-compact \(^{(3)}\) if it satisfies these properties.

\begin{enumerate}
\def\labelenumi{\alph{enumi})}
\setcounter{enumi}{1}
\item
  Let E and F be Banach spaces. Show that the set \(\mathcal{L}^{\mathrm{sc}}(\mathrm{E};\mathrm{F})\) of semi-compact mappings is a closed subspace of \(\mathcal{L}(\mathrm{E};\mathrm{F})\) .
\item
  Let E and F be Banach spaces and \(u \in \mathcal{L}(\mathrm{E};\mathrm{F})\) . Let \(E_{1}\) , \(F_{1}\) be Banach spaces, \(f \in \mathcal{L}(\mathrm{E}_{1};\mathrm{E})\) and \(g \in \mathcal{L}(\mathrm{F};\mathrm{F}_{1})\) . If \(u: E \to F\) is semi-compact, so is \(g \circ u \circ f\) .
\item
  Let E and F be Banach spaces and \(u \in \mathcal{L}(\mathrm{E};\mathrm{F})\) . If u is compact, it is semi-compact. Conversely, if E is reflexive and u is semi-compact, then u is compact.
\item
  Let E and F be Banach spaces. If the dual space E′ of E is of countable type, then every completely continuous linear map E → F is compact.
\item
  Let I be an infinite set; the identity map of \(\ell^{1}(\mathbf{I})\) is semi-compact, but is not weakly compact (cf. EVT, IV, p. 54, exerc. 15). g) The injection of \(\ell_{\mathrm{K}}^{1}(\mathbf{N})\) into \(\ell_{\mathrm{K}}^{2}(\mathbf{N})\) is semi-compact but is not compact (cf. loc. cit.).
\item
  Let E be a Banach space. Every continuous linear map \(E \rightarrow \ell_{K}^{1}(\mathbf{N})\) is semi-compact, and every continuous linear map \(\ell_{\mathrm{K}}^{1}(\mathbf{N}) \rightarrow \mathrm{E}\) is semi-compact.
(3) Some authors call them completely continuous maps.
\item
  Let E be a Banach space. If E is reflexive or if its dual is of countable type, then every continuous linear map \(\mathrm{E} \to \ell_{\mathrm{K}}^{1}(\mathbf{N})\) is compact.
\end{enumerate}

\begin{enumerate}
\def\labelenumi{\arabic{enumi})}
\setcounter{enumi}{30}
\tightlist
\item
  Let E be a Banach space.
\end{enumerate}

\begin{enumerate}
\def\labelenumi{\alph{enumi})}
\tightlist
\item
  Prove that the following conditions are equivalent:
\end{enumerate}

\begin{enumerate}
\def\labelenumi{(\roman{enumi})}
\item
  For any Banach space F, every weakly compact mapping from E into F is semi-compact.
\item
  If \((x_{n})\) is a sequence tending to 0 in E for the topology \(\sigma(\mathrm{E},\mathrm{E}^{\prime})\) and \((x_{n}^{\prime})\) a sequence tending to 0 in \(E^{\prime}\) for the topology \(\sigma(\mathrm{E}^{\prime},\mathrm{E}^{\prime\prime})\) , the sequence \(\langle x_{n},x_{n}^{\prime}\rangle\) tends to 0.
\end{enumerate}

(To prove that (i) implies (ii), one takes \(F = c_0\) . If (i) is not satisfied, there exists a weakly compact map \(u\) from E into a Banach space F and a sequence \((x_n)\) in E tending weakly to 0, such that \(\inf(\|u(x_n)\|) > 0\) ; one constructs a sequence \((y_n)\) in the unit ball of \(E'\) , convergent for \(\sigma(E', E'')\) and such that \(\langle u(x_n), y_n \rangle = \|u(x_n)\|\) , and one takes \(x_n' = y_n - \lim y_n\) .)

These conditions are satisfied in particular when \(E = \mathcal{C}(T)\) , where T is a compact space, and when \(E = L^{1}(X, \mu)\) , or \(E = L^{\infty}(X, \mu)\) , where X is a locally compact space and \(\mu\) a measure on X (INT, VI, § 2, exerc. 23). b) Prove that a Banach space possessing the above properties admits no reflexive direct factor of infinite dimension.

\begin{enumerate}
\def\labelenumi{\alph{enumi})}
\setcounter{enumi}{2}
\tightlist
\item
  Let H be a reflexive Banach space, and let B be the unit ball of H', equipped with the weak topology. Show that H can be identified with a closed non-direct subspace of \(\mathcal{C}(\mathrm{B})\) .
\end{enumerate}

\begin{enumerate}
\def\labelenumi{\arabic{enumi})}
\setcounter{enumi}{31}
\item
  A Banach space of countable type admitting a Schauder basis (EVT, IV, p. 70, exerc. 14) possesses the approximation property.
\item
  \begin{enumerate}
  \def\labelenumii{\alph{enumii})}
  \tightlist
  \item
    Prove that the strict inductive limit of a sequence of approximation spaces is an approximation space.
  \end{enumerate}
\end{enumerate}

\begin{enumerate}
\def\labelenumi{\alph{enumi})}
\setcounter{enumi}{1}
\tightlist
\item
  Prove that the locally convex direct sum of a family of approximation spaces is an approximation space.
\end{enumerate}

\begin{enumerate}
\def\labelenumi{\arabic{enumi})}
\setcounter{enumi}{33}
\tightlist
\item
  Let E be a Banach space. Prove that the following properties are equivalent:
\end{enumerate}

\begin{enumerate}
\def\labelenumi{(\roman{enumi})}
\item
  The space E' is an approximation space;
\item
  For every Banach space F, the closure of \(\mathcal{L}^{\mathrm{f}}(\mathrm{E};\mathrm{F})\) in \(\mathcal{L}(\mathrm{E};\mathrm{F})\) is \(\mathcal{L}^{\mathrm{c}}(\mathrm{E};\mathrm{F})\) ;
\item
  For every closed subspace F of \(c_{0}\) the closure of \(\mathcal{L}^{\mathrm{f}}(\mathrm{E};\mathrm{F})\) in \(\mathcal{L}(\mathrm{E};\mathrm{F})\) is \(\mathcal{L}^{\mathrm{c}}(\mathrm{E};\mathrm{F})\) .
\end{enumerate}

\begin{enumerate}
\def\labelenumi{\arabic{enumi})}
\setcounter{enumi}{34}
\tightlist
\item
  Let E be a Banach space. We say that E has the metric approximation property if the identity map of E belongs to the closure of the set of elements of norm \(\leqslant 1\) of \(\mathcal{L}^{\mathrm{f}}(\mathrm{E})\) in the space \(\mathcal{L}(\mathrm{E})\) endowed with the topology of compact convergence.
\end{enumerate}

Prove that the following properties are equivalent:

\begin{enumerate}
\def\labelenumi{(\roman{enumi})}
\item
  The space E has the metric approximation property;
\item
  The unit ball of \(\mathcal{L}(\mathrm{E})\) is the closure, for the topology of compact convergence, of the set of elements of norm \(\leqslant 1\) of \(\mathcal{L}^{\mathrm{f}}(\mathrm{E})\) ;
\item
  For every Banach space F, the unit ball of \(\mathcal{L}(\mathrm{E};\mathrm{F})\) is the closure, for the topology of compact convergence, of the set of elements of norm \(\leqslant 1\) of \(\mathcal{L}^{\mathrm{f}}(\mathrm{E};\mathrm{F})\) ;
\item
  For every Banach space F, the unit ball of \(\mathcal{L}(\mathrm{F};\mathrm{E})\) is the closure, for the topology of compact convergence, of the set of elements of norm \(\leqslant 1\) of \(\mathcal{L}^{\mathrm{f}}(\mathrm{F};\mathrm{E})\) .
\end{enumerate}

\begin{enumerate}
\def\labelenumi{\arabic{enumi})}
\setcounter{enumi}{35}
\tightlist
\item
  \begin{enumerate}
  \def\labelenumii{\alph{enumii})}
  \tightlist
  \item
    Let X be a locally compact topological space. Prove that the space \(\mathcal{C}_{0}(\mathrm{X})\) has the metric approximation property (adapt the proof of Prop. 13 of III, p. 21).
  \end{enumerate}
\end{enumerate}

\begin{enumerate}
\def\labelenumi{\alph{enumi})}
\setcounter{enumi}{1}
\tightlist
\item
  Let X be a locally compact topological space equipped with a positive measure and \(p \in [1, +\infty]\) . Prove that the space \(L^{p}(X)\) has the metric approximation property (same method).
\end{enumerate}

\begin{enumerate}
\def\labelenumi{\arabic{enumi})}
\setcounter{enumi}{36}
\tightlist
\item
  Let X be a differentiable manifold of class \(C^{r}\) where \(1 \leqslant r \leqslant \infty\) , paracompact, locally of finite dimension, and let E be a vector bundle over X of class \(C^{r}\) , locally of finite dimension. Prove that the space \(\mathcal{S}^{r}(\mathrm{X};\mathrm{E})\) has the approximation property.
\end{enumerate}

\exercisesection{}

\begin{enumerate}
\def\labelenumi{\arabic{enumi})}
\tightlist
\item
  Let \(I = [a, b]\) be a compact interval of R. For every continuous function \(f: I \to C\) and every \(y \in I\) set
\end{enumerate}

\[
\mathrm{V} (f) (y) = \int_ {a} ^ {y} f (x) d x.
\]

Show that the linear map V from \(\mathcal{C}(\mathrm{I})\) into \(\mathcal{C}(\mathrm{I})\) is compact (“Volterra operator”).

\begin{enumerate}
\def\labelenumi{\arabic{enumi})}
\setcounter{enumi}{1}
\tightlist
\item
  Let X = {]}0, 1{[} and let μ be Lebesgue measure on X. One defines a sequence \((f_{n})_{n \geqslant 1}\) of functions X → R by setting
\end{enumerate}

\[
f _ {2 n} (x) = \left\{ \begin{array}{l l} \sqrt {2} \sin (2 n \pi x) & \text { if } 0 <   x <   1 / 2 \\ - 2 ^ {n} & \text { if } 1 - 2 ^ {- 2 n} <   x <   1 - 2 ^ {- 2 n - 1} \\ 0 & \text { otherwise }, \end{array} \right.
\]

and

\[
f _ {2 n - 1} (x) = \left\{ \begin{array}{l l} \sqrt {2} \sin (2 n \pi x) & \text { if } 0 <   x <   1 / 2 \\ 2 ^ {n} & \text { if } 1 - 2 ^ {- 2 n} <   x <   1 - 2 ^ {- 2 n - 1} \\ 0 & \text { otherwise }, \end{array} \right.
\]

for all \(n \geqslant 1\) and every \(x \in X\) .

\begin{enumerate}
\def\labelenumi{\alph{enumi})}
\item
  The family \((f_n)_{n\geqslant 1}\) is orthonormal in \(\mathcal{L}^2 (\mathrm{X},\mu)\) .
\item
  For every \((x,y)\in \mathrm{X}\times \mathrm{X}\) we set
\end{enumerate}

\[
\mathrm{N} (x, y) = \sum_ {n = 1} ^ {+ \infty} 2 ^ {- n} (f _ {2 n - 1} (x) f _ {2 n - 1} (y) - f _ {2 n} (x) f _ {2 n} (y)).
\]

Show that \(|\mathrm{N}(x,y)| \leqslant 2^{3/2}\) for all \((x,y) \in \mathrm{X} \times \mathrm{X}\) . Deduce that \(\mathrm{N}_y \in \mathcal{L}^\infty(\mathrm{X})\) for all \(y \in \mathrm{X}\) and that

\[
\int_ {X} ^ {*} \| N _ {y} \| _ {\infty} ^ {2} d y <   + \infty .
\]

\begin{enumerate}
\def\labelenumi{\alph{enumi})}
\setcounter{enumi}{2}
\tightlist
\item
  Show that by setting
\end{enumerate}

\[
u _ {\mathrm{N}} (f) (y) = \int_ {\mathrm{X}} \mathrm{N} (x, y) f (x) d \mu
\]

for \(f \in \mathcal{L}^{1}(X, \mu)\) and \(y \in X\) , by passing to the quotient one defines a continuous linear map from \(\mathrm{L}^{1}(X, \mu)\) into \(\mathrm{L}^{2}(X, \mu)\) .

\begin{enumerate}
\def\labelenumi{\alph{enumi})}
\setcounter{enumi}{3}
\tightlist
\item
  The linear map \(u_{N}\) is not compact. (Consider the sequence \(g_{n}=2^{n}(f_{2n-1}-f_{2n})\) ; show that \(g_{n}\) is integrable and that the sequence \((u_{\mathrm{N}}(g_{n}))\) in \(\mathrm{L}^{2}(\mathrm{X},\mu)\) has no convergent subsequence.)
\end{enumerate}

\begin{enumerate}
\def\labelenumi{\arabic{enumi})}
\setcounter{enumi}{2}
\tightlist
\item
  Let X be a locally compact topological space, \(\mu\) a positive measure on X and \(Y \subset X\) a locally \(\mu\) -negligible subset. Let N be the characteristic function of \(X \times Y\) .
\end{enumerate}

\begin{enumerate}
\def\labelenumi{\alph{enumi})}
\tightlist
\item
  One has
\end{enumerate}

\[
\int_ {X \times X} ^ {*} | N (x, y) f (x) g (y) | d (\mu \otimes \mu) (x, y) = 0
\]

for \(f, g \in \mathrm{L}^2(\mathrm{X}, \mu)\) .

\begin{enumerate}
\def\labelenumi{\alph{enumi})}
\setcounter{enumi}{1}
\tightlist
\item
  If \(f \in \mathrm{L}^{2}(\mathrm{X}, \mu) - \mathrm{L}^{1}(\mathrm{X}, \mu)\) , the set of \(y \in X\) such that \(x \mapsto \mathrm{N}(x, y)f(x)\) is \(\mu\) -integrable coincides with Y.
\end{enumerate}

\begin{enumerate}
\def\labelenumi{\arabic{enumi})}
\setcounter{enumi}{3}
\tightlist
\item
  Let \(\alpha \in ]0,1]\) . Let I be a compact interval of \(\mathbf{R}\) . We denote by \(\mathcal{C}^{\alpha}(\mathrm{I})\) the vector subspace of \(\mathcal{C}(\mathrm{I})\) consisting of functions \(f\) such that there exists \(\mathrm{C} \geqslant 0\) satisfying
\end{enumerate}

\[
| f (x) - f (y) | \leqslant \mathrm{C} | x - y | ^ {\alpha}
\]

for all \((x,y)\in\mathrm{I}\times\mathrm{I}\) .

We equip \(\mathcal{C}^{\alpha}(\mathrm{I})\) with the norm

\[
\| f\|_{\alpha} = \sup_{x\in \mathrm{I}}|f(x)| + \sup_{\substack{(x,y)\in \mathrm{I}\times \mathrm{I}\\ x\neq y}}\frac{|f(x) - f(y)|}{|x - y|^{\alpha}}.
\]

\begin{enumerate}
\def\labelenumi{\alph{enumi})}
\item
  The space \(\mathcal{C}^{\alpha}(\mathrm{I})\) is a Banach space.
\item
  The canonical injection \(\mathcal{C}^{\alpha}(\mathrm{I})\to \mathcal{C}(\mathrm{I})\) is compact.
\item
  If \(0 < \alpha < \beta \leqslant 1\) then the canonical injection \(\mathcal{C}^{\beta}(\mathrm{I}) \to \mathcal{C}^{\alpha}(\mathrm{I})\) is compact.\\
\item
  The Sobolev space \(\mathrm{H}^{1}(\mathring{\mathrm{I}})\) (cf. IV, p. 221, no. 14) is identified with a subspace of \(\mathcal{C}^{1/2}(\mathrm{I})\) but the canonical injection \(\mathrm{H}^{1}(\mathring{\mathrm{I}}) \to \mathcal{C}^{1/2}(\mathrm{I})\) is not compact.
\end{enumerate}

\exercisesection{}

\begin{enumerate}
\def\labelenumi{\arabic{enumi})}
\item
  Let E and F be Hilbert spaces and let u be a continuous linear map from E into F. Show that u is a Fredholm map if and only if the image of u is closed and the kernels of u and \(u^{*}\) are finite-dimensional.
\item
  Let \(p\) be a real number such that \(1 \leqslant p < +\infty\) . Let \(u\) be a diagonal endomorphism of \(\ell_{\mathrm{K}}^{p}(\mathbf{N})\) . Show that \(u\) is a Fredholm endomorphism if and only if it is a Riesz endomorphism. In particular, every diagonal Fredholm endomorphism of \(\ell_{\mathrm{K}}^{p}(\mathbf{N})\) has index zero.
\item
  \begin{enumerate}
  \def\labelenumii{\alph{enumii})}
  \tightlist
  \item
    There exists a Banach space E, a continuous automorphism u of E, and a Riesz endomorphism v of E such that \(u \circ v\) is not a Riesz endomorphism. There even exist such examples for which \(u \circ v - v \circ u\) has finite rank.
  \end{enumerate}
\end{enumerate}

\begin{enumerate}
\def\labelenumi{\alph{enumi})}
\setcounter{enumi}{1}
\tightlist
\item
  There exists a Banach space E, an automorphism u of E, and a finite-rank endomorphism h such that \(u + h\) is not a Riesz endomorphism.
\end{enumerate}

\begin{enumerate}
\def\labelenumi{\arabic{enumi})}
\setcounter{enumi}{3}
\tightlist
\item
  Let E be an infinite-dimensional complex Hilbert space.
\end{enumerate}

\begin{enumerate}
\def\labelenumi{\alph{enumi})}
\item
  The group \(\mathbf{U}(\mathrm{E})\) of unitary operators on E, equipped with the topology induced by the norm of \(\mathcal{L}(\mathrm{E})\) , is path-connected. (Show that if u is unitary, there exists a Hermitian operator v such that \(u = \exp(iv)\) .)
\item
  The group of linear isomorphisms of E is path-connected. (Use the polar decomposition, cf. I, p. 139, def. 3.)
\item
  For every \(n \in Z\) the set of Fredholm endomorphisms of E of index n is open and connected. It is nonempty.
\end{enumerate}

\begin{enumerate}
\def\labelenumi{\arabic{enumi})}
\setcounter{enumi}{4}
\tightlist
\item
  Let E and F be infinite-dimensional Banach spaces. A continuous linear map \(u: \mathrm{E} \to \mathrm{F}\) is said to be strictly singular if there does not exist a closed subspace \(\mathrm{F}_1 \subset \mathrm{E}\) of infinite dimension such that \(u\) induces, by passage to subspaces, an isomorphism from \(\mathrm{F}_1\) onto \(u(\mathrm{F}_1)\) .
\end{enumerate}

\begin{enumerate}
\def\labelenumi{\alph{enumi})}
\tightlist
\item
  Let u be a continuous linear map from E into F. The following conditions are equivalent:
\end{enumerate}

\begin{enumerate}
\def\labelenumi{(\roman{enumi})}
\item
  The map \(u\) is strictly singular;
\item
  For every closed subspace \(F_{1}\) of E of infinite dimension, there does not exist a real number \(\delta > 0\) such that \(\|u(x)\| \geqslant \delta\|x\|\) for all \(x \in F_{1}\) .
\end{enumerate}

\begin{enumerate}
\def\labelenumi{\alph{enumi})}
\setcounter{enumi}{1}
\item
  Let \(u \colon \mathrm{E} \to \mathrm{F}\) be a continuous linear map such that the image of \(u\) is closed and the kernel of \(u\) is finite-dimensional. If \(v \colon \mathrm{E} \to \mathrm{F}\) is strictly singular, then the image of \(u + v\) is closed and its kernel is finite-dimensional.
\item
  The set of strictly singular endomorphisms of E is a closed two-sided ideal of \(\mathcal{L}(\mathrm{E})\) .
\item
  Every compact linear map is strictly singular. If E and F are Hilbert spaces, then every strictly singular map \(E \rightarrow F\) is compact.
\item
  There exist Banach spaces E and F and a strictly singular non-compact map \(E \rightarrow F\) .
\end{enumerate}

\(f)\) Let \(p\) and \(r\) be real numbers such that \(1 \leqslant p < r\) . Every continuous linear map \(\ell_{\mathrm{K}}^{r}(\mathbf{N}) \to \ell_{\mathrm{K}}^{p}(\mathbf{N})\) is strictly singular. (Use exercise 10 of III, p. 106.)

\begin{enumerate}
\def\labelenumi{\arabic{enumi})}
\setcounter{enumi}{5}
\item
  Let E be a complex Hilbert space. Let u be a Fredholm endomorphism of E. If u is normal, then \(\text{ind}(u) = 0\) .
\item
  Let E and F be Banach spaces, and let u be an element of \(\mathcal{L}(\mathrm{E};\mathrm{F})\) . Prove that u is a Fredholm map if and only if the following conditions are satisfied:
\end{enumerate}

\begin{enumerate}
\def\labelenumi{(\roman{enumi})}
\tightlist
\item
  There exists a continuous seminorm p on E and \(C \geqslant 0\) such that the canonical map from E into the separated completion of \((E, p)\) is compact, and
\end{enumerate}

\[
\| x \| \leqslant \mathrm{C} \| u (x) \| + p (x)
\]

for all \(x \in E\) ;

\begin{enumerate}
\def\labelenumi{(\roman{enumi})}
\setcounter{enumi}{1}
\tightlist
\item
  There exists a continuous seminorm \(q\) on \(\mathrm{F}'\) and \(\mathrm{C}' \geqslant 0\) such that the canonical map from \(\mathrm{F}'\) into the separated completion of \((\mathrm{F}', q)\) is compact, and
\end{enumerate}

\[
\| y ^ {\prime} \| \leqslant \mathrm{C} ^ {\prime} \| ^ {t} u (y ^ {\prime}) \| + q (y ^ {\prime})
\]

for all \(y' \in F'\) .

\begin{enumerate}
\def\labelenumi{\arabic{enumi})}
\setcounter{enumi}{7}
\tightlist
\item
  Let U be a connected open subset of C. We denote by \(\mathcal{H}(\mathrm{U})\) the space of holomorphic functions on U equipped with the topology of compact convergence (EVT III, p. 10, example \(d\) ).
\end{enumerate}

\begin{enumerate}
\def\labelenumi{\alph{enumi})}
\item
  Let \(a \in \mathcal{H}(U)\) . Multiplication by a is a Fredholm endomorphism of \(\mathcal{H}(U)\) if and only if the set of zeros of a is finite. Then compute its index.
\item
  Let \(\mathcal{U} = (\mathrm{U}_{i})_{i \in \mathrm{I}}\) a covering of U by nonempty connected open sets. We set
\end{enumerate}

\[
\begin{array}{l} \mathrm{I} _ {2} = \{(i, j) \in \mathrm{I} \times \mathrm{I} | \mathrm{U} _ {i} \cap \mathrm{U} _ {j} \neq \varnothing \}, \\ \mathrm{I} _ {3} = \{(i, j, k) \in \mathrm{I} \times \mathrm{I} \times \mathrm{I} | \mathrm{U} _ {i} \cap \mathrm{U} _ {j} \cap \mathrm{U} _ {k} \neq \varnothing \}. \end{array}
\]

We consider the vector spaces

\[
\begin{array}{l} \mathrm{Z} ^ {1} (\mathcal {U}; \mathbf {C}) = \{(c _ {i, j}) \in \mathbf {C} ^ {\mathrm{I} _ {2}} | c _ {i j} + c _ {j k} + c _ {k i} = 0 \text { for all } (i, j, k) \in \mathrm{I} _ {3} \} \\ \mathrm{B} ^ {1} (\mathcal {U}; \mathbf {C}) = \{(c _ {i, j}) \in \mathbf {C} ^ {\mathrm{I} _ {2}} | \text { there exists } c ^ {\prime} \in \mathbf {C} ^ {\mathrm{I}} \text { such that } \end{array}
\]

\[
c _ {i j} = c _ {i} ^ {\prime} - c _ {j} ^ {\prime} \text {   for all   } (i, j) \in \mathrm{I} _ {2} \}.
\]

One has \(\mathrm{B}^1 (\mathcal{U};\mathbf{C})\subset \mathrm{Z}^1 (\mathcal{U};\mathbf{C})\) . We set

\[
\mathrm{H} ^ {1} (\mathcal {U}; \mathbf {C}) = \mathrm{Z} ^ {1} (\mathcal {U}; \mathbf {C}) / \mathrm{B} ^ {1} (\mathcal {U}; \mathbf {C}).
\]

Let \(\mathcal{V}=(\mathrm{V}_{\alpha})_{\alpha\in\mathrm{A}}\) be a covering of U by nonempty connected open sets, finer than U, and let \(\varrho\colon A\to I\) be a map such that \(V_{\alpha}\subset U_{\varrho(\alpha)}\) for all \(\alpha\in A\) .

The map \((c_{ij})_{i,j} \mapsto (c_{\varrho(\alpha)\varrho(\beta)})_{\alpha,\beta}\) induces a linear map, independent of the choice of \(\varrho\) , from \(\mathrm{H}^1(\mathcal{U};\mathbf{C})\) into \(\mathrm{H}^1(\mathcal{V};\mathbf{C})\) .

We denote by \(H^{1}(U; \mathbf{C})\) the inductive limit of vector spaces \(H^{1}(\mathcal{U}; \mathbf{C})\) equipped with these linear maps, as U ranges over the set of coverings of U by nonempty connected open sets.

\begin{enumerate}
\def\labelenumi{\alph{enumi})}
\setcounter{enumi}{2}
\tightlist
\item
  Let D be the endomorphism of \(\mathcal{H}(\mathrm{U})\) defined by \(Df = f'\) . Show that there exists an isomorphism of vector spaces from the cokernel of D onto \(\mathrm{H}^1 (\mathrm{U};\mathbf{C})\) . Deduce from this that D is a Fredholm endomorphism if and only if the space \(\mathrm{H}^1 (\mathrm{U};\mathbf{C})\) is finite-dimensional. Then compute the index of D.
\end{enumerate}

\exercisesection{}

\begin{enumerate}
\def\labelenumi{\arabic{enumi})}
\item
  Let E and F be Banach spaces and u a continuous linear map \(E \rightarrow F\) . We have \(((u)) > 0\) if and only if the image of u is closed in F.
\item
  Let \(p \in [1, +\infty]\) , and \(\mathrm{E} = \ell_{\mathbf{C}}^{p}(\mathbf{N})\) . Let u be a diagonal endomorphism of E. For every \(\lambda \in C\) , compute \(((u - \lambda 1_{\mathrm{E}}))\) in terms of the spectrum of u.
\item
  Let \(p \in [1, +\infty]\) and \(\mathrm{E} = \ell_{\mathrm{K}}^{p}(\mathbf{N})\) and let \(u\) be the shift endomorphism
\end{enumerate}

\[
(x _ {0}, x _ {1}, \dots ,) \mapsto (0, x _ {0}, x _ {1}, \dots)
\]

of E. Show that u is a Fredholm endomorphism of index 1 and that \(((u)) = 1\) .

\begin{enumerate}
\def\labelenumi{\arabic{enumi})}
\setcounter{enumi}{3}
\item
  Let E be a Banach space. The map \(u \mapsto ((u))\) of \(\mathcal{L}(\mathrm{E})\) into R is continuous if and only if E has dimension \(\leqslant 1\) .
\item
  Let E be a Hilbert space and \(u \in \mathcal{L}(\mathrm{E})\) . We have \(((u)) = \inf(\mathrm{Sp}(|u|) - \{0\})\) (cf. I, p. 139, def. 3 for the definition of \(|u|\) ).
\item
  Let \(E = \ell^{2}(\mathbf{N})\) and \((e_{n})_{n \in \mathbf{N}}\) the canonical orthonormal basis of E. Let u be the diagonal endomorphism of E such that \(u(e_{0}) = e_{0}\) and \(u(e_{n}) = 2e_{n}\) if \(n \geqslant 1\) . Show that \(((u)) = 1\) and that \(u + h\) is a Fredholm endomorphism for every h such that \(\|h\| < 2\) .
\item
  Let E and F be Banach spaces and let u be an element of \(\mathcal{Q}\mathcal{M}(\mathrm{E};\mathrm{F})\) (resp. of \(\mathcal{Q}\mathcal{E}(\mathrm{E};\mathrm{F})\) ). Let h be an element of \(\mathcal{L}(\mathrm{E};\mathrm{F})\) . Suppose there exist positive real numbers a and b such that \(a + b((u))^{-1} < 1\) so that
\end{enumerate}

\[
\| h (x) \| \leqslant a \| u (x) \| + b \| x \|\tag{2}
\]

for all \(x \in E\) .

Prove that \(u + h \in \mathcal{Q}\mathcal{M}(\mathrm{E};\mathrm{F})\) (resp. \(u + h \in \mathcal{Q}\mathcal{E}(\mathrm{E};\mathrm{F})\) ) and

\[
\begin{array}{c} \dim \operatorname{Ker} (u + h) \leqslant \dim \operatorname{Ker} (u), \\ \dim \operatorname{Coker} (u + h) \leqslant \dim \operatorname{Coker} (u), \\ \operatorname{ind} (u + h) = \operatorname{ind} (u). \end{array}
\]

(Reduce to the case where \(a\) and \(b\) are \(>0\) and let \(c > 0\) be such that \(a + b/c < 1\) ; consider the norm \(p\) on E defined by

\[
p (x) = a \| u (x) \| + b \| x \| \quad \text {   for   } x \in \mathrm{E},
\]

and \(E_{1}\) the Banach space E equipped with the norm p; we have \(\mathcal{Q}\mathcal{M}(E_{1};F)=\mathcal{Q}\mathcal{M}(E;F)\) , \(\mathcal{Q}\mathcal{E}(E_{1};F)=\mathcal{Q}\mathcal{E}(E;F)\) and \(\mathcal{L}(E_{1};F)=\mathcal{L}(E;F)\) . Then apply lemma 2 of III, p. 62.)

\begin{enumerate}
\def\labelenumi{\arabic{enumi})}
\setcounter{enumi}{7}
\item
  Let E be a Hilbert space. Prove directly Theorem 4 of III, p. 64 without using the Borsuk--Ulam theorem.
\item
  Let \(r \geqslant 1\) be an integer. If \(f: \mathrm{M}_1 \to \mathrm{M}_2\) is a morphism of class \(\mathrm{C}^r\) of manifolds of class \(\mathrm{C}^r\) , one says that \(x \in \mathrm{M}_1\) is a critical point of \(f\) if \(f\) is not a submersion at \(x\) (VAR, R1, 5.9.1). Let X be the set of critical points of \(f\) . The set \(f(\mathrm{X})\) in \(\mathrm{M}_2\) is called the set of critical values of \(f\) , and \(f(\mathrm{M}_1) = f(\mathrm{X})\) is called the set of regular values of \(f\) .
\end{enumerate}

If \(M_{1}\) and \(M_{2}\) are real manifolds of class \(C^{r}\) (VAR, R1, 5.1.5, p. 35), a morphism of manifolds \(f: M_{1} \to M_{2}\) , of class \(C^{r}\) , is said to be Fredholm if, for every \(x \in M_{1}\) , the tangent linear map \(\mathrm{T}_{x}(f)\) is a Fredholm endomorphism of the Banach space \(\mathrm{T}_{x}(\mathrm{M}_{1})\) into the Banach space \(\mathrm{T}_{f(x)}(\mathrm{M}_{2})\) . a) Let \(f: M_{1} \to M_{2}\) be a Fredholm morphism. If \(M_{1}\) is connected, the index \(\operatorname{ind}(\mathrm{T}_{x}(f))\) is independent of \(x \in M_{1}\) . It is called the index of the Fredholm morphism f and is denoted \(\operatorname{ind}(f)\) .

\begin{enumerate}
\def\labelenumi{\alph{enumi})}
\setcounter{enumi}{1}
\tightlist
\item
  Let U be an open subset of a Banach space E such that \(0 \in U\) and let F be a Banach space. Let \(f: U \to F\) be a Fredholm morphism of class \(C^{r}\) and let \(u = \mathrm{T}_{0}(f)\) . There exist a Banach space \(E_{1}\) , a chart \((\mathrm{U}_{1}, \varphi, \mathrm{E}_{1} \times \mathrm{Ker}(u))\) of U such that \(0 \in U_{1}\) , a chart \((\mathrm{U}_{2}, \psi, \mathrm{E}_{1} \times \mathrm{Coker}(u))\) of F such that \(f(0) \in U_{2}\) and a Fredholm morphism \(g: \varphi(\mathrm{U}_{1}) \to \mathrm{E}_{1} \times \mathrm{Coker}(u)\) of class \(C^{r}\) such that
\end{enumerate}

\[
(\psi \circ f \circ \varphi^ {- 1}) (a, b) = (a, g (a, b))
\]

for all \((a,b)\in \varphi (\mathrm{U}_1)\subset \mathrm{E}_1\times \mathrm{Ker}(u).\)

\begin{enumerate}
\def\labelenumi{\alph{enumi})}
\setcounter{enumi}{2}
\item
  Let \(f \colon \mathrm{M}_1 \to \mathrm{M}_2\) be a Fredholm morphism of class \(\mathbf{C}^r\) between manifolds of class \(\mathbf{C}^r\) . For every \(y \in \mathbf{M}_2\) , there exists an open neighborhood \(\mathrm{V}\) of \(y\) such that \(f\) defines, by passage to subspaces, a proper map \(f^{-1}(\mathrm{V}) \to \mathrm{V}\) .
\item
  Let U be an open subset of a Banach space E containing 0 and let F be a Banach space. Let \(f \colon \mathrm{U} \to \mathrm{F}\) be a Fredholm morphism of class \(\mathrm{C}^{\infty}\) . For every \(x \in \mathrm{U}\) and every neighborhood V of \(f(x)\) in F, there exists a regular value of \(f\) in V. (Use the second Sard theorem, VAR, R2, 10.1.3, p. 37.) e) Let \(f \colon \mathrm{M}_1 \to \mathrm{M}_2\) be a Fredholm morphism of class \(\mathrm{C}^{\infty}\) of manifolds of class \(\mathrm{C}^{\infty}\) . If \(\mathrm{M}_1\) and \(\mathrm{M}_2\) admit countable bases, then the set of critical values of \(f\) is meager ("Smale's theorem").
\end{enumerate}

\(f)\) Let \(f\colon \mathrm{M}_1\to \mathrm{M}_2\) be a Fredholm morphism of class \(\mathrm{C}^{\infty}\) of manifolds of class \(\mathrm{C}^{\infty}\) admitting countable bases. If \(\operatorname {ind}(f) < 0\) , then the image of \(f\) has empty interior.

\(g)\) Let \(f\colon \mathrm{M}_1\to \mathrm{M}_2\) be a Fredholm morphism of class \(\mathrm{C}^{\infty}\) of manifolds of class \(\mathrm{C}^{\infty}\) admitting countable bases. For every \(y\in \mathrm{M}_2\) outside a meager set, the set \(f^{-1}(y)\) is a pure submanifold of dimension \(\operatorname {ind}(f)\) .